\documentclass[../ppgc-tese.tex]{subfiles}
%\graphicspath{{\subfix{../images/}}}
\begin{document}

% =============================================================================
% Chapter source: "Activation-Space Adversarial Detection in a Vision
% Transformer Measures Displacement, Not Meaning" (TMLR submission),
% repository interp_adv_tmlr_2026, main.tex at commit 833c357 (2026-11-16).
% To resynchronize: git -C ../interp_adv_tmlr_2026 diff 833c357 -- main.tex figures/
%
% This file was rebuilt on that source on 2026-11-16 (Phase 1). It replaces the
% chapter written from the July version of the paper (v06), which is in the git
% history of this repository. What changed with respect to that version:
%   - claim: the probe reads the displacement ALONG A LEARNED DIRECTION; a large
%     displacement is sufficient for detection, not necessary; "meaning" is
%     tested as the change of the label;
%   - structure: five steps, each opened by an objection (was: twelve objections
%     and a three-degree specificity ladder, both removed);
%   - notation: \Delta a (was \Delta x); vocabulary: input-matched and
%     displacement-matched control, detection advantage, watermark control,
%     head split, pre-specified (was pre-registered);
%   - new experiments: APGD-CE replication and the adaptive APGD attack.
%
% Phase 2 (2026-11-16) put back the material that only the thesis has, revised
% to the new claim and vocabulary: the verified input match per family
% (tab:adv-input-match), the sweeps over layers and representations
% (sec:adv-layers), the pre-specified rules at the frozen budget
% (sec:adv-provenance) and the two measurement bugs (sec:adv-bugs).
% Dropped from the July version, with the reason: the ViT-L/16 scale check and
% the TopK SAE (old protocol; the paper removed both), the note on a non-linear
% probe recovering 0.14 AUC (earlier protocol, not re-run), fig_gap_regimes
% (superseded by fig_controls), fig_jspace (three numbers that are in the text)
% and fig_ladder (the ladder was cut).
%
% Thesis notation that differs from the paper: the maximum in \dauc runs over
% benign families b (the paper uses f, which is the sub-sampling fraction of
% Chapter 3); N is the number of examples and K the number of folds.
%
% Evidence-status legend, kept in comments next to each claim:
%   [EST]  established by data on disk in the research repository
%   [PRE]  pre-specified hypothesis or decision rule (never relaxed)
%   [EXPL] exploratory run whose expected outcome was written before execution
%   [EXT]  external cited claim
% Every number is copied from main.tex; its \src trailer is kept as "% src:".
%
% Text that depends on results still open in the paper is wrapped in
% \pendingpaper{<action point>}{...} (defined in ppgc-tese.tex; prints its
% second argument unchanged). Open points: A1 (APGD-CE below the ceiling),
% A3 (attack-vs-benign probe), A4 (benign control matched in input L_inf to the
% minimal-displacement attack), B13 (single-detector result as the main result).
% =============================================================================

Sparse autoencoders (SAEs) promise interpretable features in vision models, and
one proposed application is the detection of adversarial attacks from the
activations of a model. This chapter asks not whether such a detector can be
built, but what it measures. We train a linear probe on the mean-pooled
activations of a ViT-B/16 classifier. The usual control matches benign
perturbations to the attack in input norm, and with this control the attack is
more detectable than every benign perturbation. This control is not matched in
what the probe reads: at equal input norm, the attack moves the activations at
the probed layer $3$ to $10\times$ more than a benign perturbation does. When
the benign perturbations are matched to the attack in this activation
displacement, each benign family is at least as separable from clean images as
the attack (Figure~\ref{fig:adv-controls}). The probe detects successful and
failed attacks equally, and the SAE is not necessary for any of these results.
% [EST — preview of sec:adv-usual, sec:adv-amplification, sec:adv-main, sec:adv-reads]

The size of the displacement does not determine detection. A smaller
displacement does not evade the probe: an attack built to minimize the
displacement reaches the benign scale, and the probe still detects it. A
different direction does evade it: an attacker who knows the probe evades it at
the budget of the reference attack, but the displacement stays at the attack
scale. In this setting, activation-space adversarial detection reads the
displacement along a learned direction, not adversarial meaning, which we test
as the change of the label.
% [EST — preview of sec:adv-attackers]

\begin{figure}[ht]
  \centering
  \includegraphics[width=0.7\linewidth]{figures/interp_adv/fig_controls.pdf}
  \caption{Detection AUC against mean $\lVert\Delta a\rVert$ at
  \texttt{mlp\_10}, probe on the SAE latents $z$, for the benign families in the
  input-matched (open) and displacement-matched (filled) controls, and for FGSM
  ($\varepsilon{=}0.004$). Within each benign family, the AUC rises with the
  displacement, and at equal mean displacement every benign family is at least
  as detectable as FGSM. The AUCs and the input-matched $\dauc$ are from one
  seed (over five seeds, $0.204\pm0.014$); the displacement-matched $\dauc$ is
  the mean $\pm$ standard deviation over five seeds.}
  % src: figures/src/fig_controls.py (paper repository)
  \label{fig:adv-controls}
\end{figure}

Section~\ref{sec:adv-hypotheses} states the confound, the two hypotheses and
the structure of the argument. Section~\ref{sec:adv-background} gives the
background specific to this study, and Section~\ref{sec:adv-instrument} defines
the model, the displacement, the perturbation families, the two matched
controls and the detection protocol. Sections~\ref{sec:adv-usual}
to~\ref{sec:adv-attackers} give the results in five steps, and each step opens
with the objection of a skeptical reader that it answers.
Section~\ref{sec:adv-layers} gives the sweeps over layers and representations,
Section~\ref{sec:adv-provenance} gives the provenance of the tests and the
pre-specified rules, Section~\ref{sec:adv-bugs} documents two measurement
bugs, Section~\ref{sec:adv-reconciliation} relates the results to deployment
work,
and Section~\ref{sec:adv-closing} states the limitations and the reading for
the thesis.

% =============================================================================
\section{Two hypotheses and a confound}
\label{sec:adv-hypotheses}

Adversarial examples remain among the most studied failure modes of deep vision
models: small input perturbations that flip a confident prediction
\cite{goodfellow2015fgsm,costa2024adversarial_survey}. A recurring hope is that
the model's own internal representations betray the attack: that somewhere in
activation space there is a signal for ``this input is adversarial'', and that
SAEs, by decomposing activations into interpretable features, can expose it
\cite{rao2024discover_then_name,pach2025sae_monosemantic_vlm}. On this view,
adversarial detection becomes a form of introspection: the model, read through
its features, ``knows'' it is under attack. % [EXT/context]

\paragraph*{The confound.}
An adversarial example is a perturbation of a clean image, and it changes the
output of the model. A detector can read either property, or both. Usual tests
compare clean inputs with adversarial inputs only, which differ in the two
properties at the same time, so they cannot separate them
\cite{metzen2017detecting_adversarial,wang2026saegis}; some tests add benign
noise matched to the attack in input norm
\cite{feinman2017adversarial_artifacts,ma2018lid_adversarial}. To separate the
two properties, we compare the attack with benign perturbations of the same
magnitude. The question is thus not whether such a detector can be built, but
what it measures, which is a question of construct validity
\cite{cronbach1955construct_validity,jacobs2021measurement_fairness}.

\paragraph*{Two definitions of ``same magnitude''.}
The first is the \emph{input norm}, the size of the change to the image. The
second is the \emph{activation displacement} $\lVert\Delta a\rVert$, the size
of the change in the activation at the layer that the detector reads
(Equation~\eqref{eq:adv-dx}). A linear probe on mean-pooled activations
responds only to $\Delta a$ at its layer, and the two definitions are not
equivalent: at equal input $L_2$, FGSM moves the activation at
\texttt{mlp\_10} by $2.00$, and the four benign families move it by
$0.19$--$0.59$ (Section~\ref{sec:adv-amplification}). We therefore use two
controls: the \emph{input-matched control} matches the input $L_2$, as usual,
and the \emph{displacement-matched control} matches $\lVert\Delta a\rVert$ at
the probed layer.

\subsection{The two hypotheses}
\label{sec:adv-two-hypotheses}

Two hypotheses compete for the same observations.

\paragraph*{The introspection hypothesis ($H_{M}$).}
The probe reads adversarial \emph{meaning}: the property that makes an example
adversarial, which we operationalize as the change of the label. Under $H_{M}$,
the attack stays more detectable than benign perturbations that are equal in
the quantity that the probe reads, successful attacks are more detectable than
failed ones, and the signal is in the part of the displacement that changes the
logits. This is the hypothesis that was written first: before any run, two
predictions of $H_{M}$ were pre-specified with fixed decision rules, H1
(detection does not depend on the attack budget) and H2 (a large detection
advantage in the input-matched control). Section~\ref{sec:adv-provenance}
gives the rules and their outcome. % [PRE]

\paragraph*{The displacement hypothesis ($H_{D}$).}
The probe reads the activation displacement $\Delta a$ along a direction that
it learns, and the change of the label has no role in what it reads.
Adversarial optimization is one cause of a large displacement, and nothing
more. $H_{D}$ has two parts, and each can fail. First, at equal
$\lVert\Delta a\rVert$ the attack has no detection advantage over benign
perturbations. Second, detection does not depend on the functional effect of
the attack. The scope of the instrument, a linear probe on mean-pooled
activations, is part of the statement: $H_{D}$ makes no claim about what a
non-linear or a token-level probe would see (Section~\ref{sec:adv-scope}).

\paragraph*{What $H_{D}$ does not say.}
$H_{D}$ does not say that the detection AUC is a function of
$\lVert\Delta a\rVert$ alone. An earlier formulation of this study did say so,
and the data reject it: the scalar norms are near chance
(Section~\ref{sec:adv-reads}), an attack with a displacement at the benign
scale is detected, and an attack with a displacement at the attack scale evades
a fixed probe (Section~\ref{sec:adv-attackers}). A large displacement is
\emph{sufficient} for detection by a probe trained on that family: every
perturbation at the attack scale that we tested is detectable. It is not
\emph{necessary}. The chapter also does not claim that attacks and benign
perturbations are indistinguishable in activation space;
\pendingpaper{A3}{we did not train an attack-vs-benign probe, and the collapse
test of Section~\ref{sec:adv-reads} suggests that the family of a perturbation
can be decoded.}

\subsection{The structure of the argument}
\label{sec:adv-flow}

A detector has discriminant validity for ``adversarial'' only if it separates
attacks from benign perturbations that are equal in the quantity that it reads
\cite{campbell1959convergent_discriminant_validation}. Each of the five steps
states one objection and gives an experiment whose outcome would differ if the
objection were correct
\cite{platt1964strong_inference,shadish2002experimental_designs}
(Table~\ref{tab:adv-flow}). The order of the steps is the order of the
argument, not the order in which the experiments were run; the tests of $H_{D}$
are exploratory, and Section~\ref{sec:adv-provenance} gives the status of each
test.

\begin{table}[ht]
  \centering
  \caption{Structure of the results, in five steps. Each step opens with the
  objection of a skeptical reader that it answers.}
  \label{tab:adv-flow}
  \small
  \begin{tabular}{@{}>{\raggedright\arraybackslash}p{0.28\linewidth}>{\raggedright\arraybackslash}p{0.22\linewidth}>{\raggedright\arraybackslash}p{0.42\linewidth}@{}}
    \toprule
    Step and objection & Experiment & Answer \\
    \midrule
    1 (\S\ref{sec:adv-usual}). ``The input-norm match is the standard control, and it shows a detection advantage.'' & Input-matched control, for FGSM, PGD and APGD-CE & The usual control reports $\dauc=0.199$ on $z$ for FGSM, and a positive $\dauc$ for each attack. \\
    \addlinespace
    2 (\S\ref{sec:adv-amplification}). ``The benign perturbations have the same input norm, so they have the same magnitude.'' & $\lVert\Delta a\rVert$ of each family; CW-$L_2$; PGD budget grid & At equal input $L_2$ the attack has $3$--$10\times$ the $\lVert\Delta a\rVert$ of the benign families. \\
    \addlinespace
    3 (\S\ref{sec:adv-main}). ``The displacement match is the wrong control.'' ``The result is a property of the SAE or of the pooling.'' & Displacement-matched control; representations; token-level probe & Benign families matched to FGSM in $\lVert\Delta a\rVert$ at \texttt{mlp\_10} are at least as detectable as FGSM, in five seeds. The SAE is not needed, and pooling hides only localized signal. \\
    \addlinespace
    4 (\S\ref{sec:adv-reads}). ``Matching $\lVert\Delta a\rVert$ removes the adversarial signal itself.'' ``$\dauc<0$, so something remains.'' & Success split; head split; scalar norms and Mahalanobis distance; collapse test & Successful and failed attacks are equally detectable, and the probe does not need the component that changes the logits. What remains is a direction: Gaussian noise scores as high as FGSM, and blur and JPEG score low. \\
    \addlinespace
    5 (\S\ref{sec:adv-attackers}). ``An attacker can evade the probe with a low-displacement example.'' & Minimal-displacement attack; adaptive APGD & An attack at benign-scale $\lVert\Delta a\rVert$ is still detected; an attacker who knows the probe evades it by direction, at attack-scale $\lVert\Delta a\rVert$. \emph{Open:} the input-$L_\infty$-matched benign control. \\
    \bottomrule
  \end{tabular}
\end{table}

% NOTE (2026-11-16): the subsection "Relation to the thesis question" of the
% July version was removed. It quoted the diagnostics research question of the
% proposal-era introduction ("under pre-registered tests, when do sparse-feature
% signatures isolate a structurally specific perturbation..."), which is to be
% rewritten (Phase 4). The reading for the thesis is in sec:adv-closing.

% =============================================================================
\section{Background specific to this study}
\label{sec:adv-background}
% From the Related Work of the paper. SAE fundamentals are in the Background
% chapter; only what that chapter lacks is kept here. All paraphrase; no
% verbatim sentences from cited work.

\paragraph*{Detecting adversarial examples in representation space.}
A small crafted perturbation can change the prediction of a model
\cite{goodfellow2015fgsm,madry2018pgd,carlini2017cw_attack}. Many detectors
read the attack from internal statistics: auxiliary classifiers
\cite{metzen2017detecting_adversarial}, density and uncertainty estimates
\cite{feinman2017adversarial_artifacts}, input squeezing
\cite{xu2018feature_squeezing}, local intrinsic dimensionality
\cite{ma2018lid_adversarial}, reconstruction \cite{meng2017magnet}, and
class-conditional Mahalanobis distances \cite{lee2018mahalanobis}. The last is
the baseline of Section~\ref{sec:adv-reads}; it detects out-of-distribution
inputs and attacks with one score, an early sign that such detectors respond to
a generic departure from clean data. Most of these detectors fail against
adaptive attackers
\cite{carlini2017bypassing_detection,sharma2018bypassing_feature_squeezing,athalye2018obfuscated_gradients,tramer2020adaptive_attacks},
detection is about as hard as robust classification
\cite{tramer2022detection_hardness}, and guidelines for rigorous evaluation
followed \cite{carlini2019evaluating_robustness}. Work on ViT robustness acts on
the architecture
\cite{bhojanapalli2021vit_robustness,naseer2021vit_properties,lovisotto2022attention_harmful,huang2023patchcensor}.
This literature rarely asks what a working detector reads. Some detectors scale
a noise baseline to the input norm of the attack
\cite{feinman2017adversarial_artifacts,ma2018lid_adversarial}, but the match is
not verified for each attack family, and no control matches the activation
displacement, so displacement and adversarial effect stay confounded. % [EXT]

\paragraph*{Controls for what a probe reads.}
High probe accuracy does not show that a representation encodes a property: a
probe needs a control that removes the property and keeps everything else
\cite{hewitt2019control_tasks,belinkov2022probing_classifiers}, and an
explanation method needs a test that it can fail
\cite{adebayo2018sanity_checks_saliency}. In SAE probing, simple baselines on
the raw activations match SAE probes \cite{kantamneni2025sae_probing}, which
agrees with the comparison of $a$ and $z$ in Section~\ref{sec:adv-scope}. In
out-of-distribution detection, feature-based scores answer a different question
from the one they are used for: a far feature-space distance is not the same as
being out of distribution \cite{li2025ood_wrong_questions}, and the feature
norm behaves like the confidence of a hidden classifier
\cite{park2023feature_norm_ood}. The benign perturbations matched in
displacement are the control task of an adversarial-detection probe. % [EXT]

\paragraph*{Sparse autoencoders and adversarial inputs.}
The Background chapter (Chapter~\ref{chpt:background}) introduces SAEs as
overcomplete dictionaries of sparse features. On vision backbones the methods
have matured: concept discovery and naming
\cite{rao2024discover_then_name,gorton2024curve_detectors_sae,pach2025sae_monosemantic_vlm,lim2025sae_visual_concepts,zaigrajew2025clip_hierarchical_sae,stevens2025testable_vision_sae},
dictionary stability and alignment
\cite{fel2025archetypal_sae,thasarathan2025universal_sae}, and steering of ViT
features \cite{joseph2025steering_clip_sae,han2025causal_sae_vision}. This
chapter uses the SAE as a measuring instrument. SAEgis \cite{wang2026saegis}
is, to our knowledge, the only SAE-based adversarial detector
(Section~\ref{sec:adv-reconciliation}). Adjacent work removes unsafe concepts
with SAEs \cite{geng2025sae_unsafe_concepts}, attacks SAE concept
representations \cite{li2026sae_concept_robustness}, and links superposition to
adversarial vulnerability \cite{bereska2025superposition_lossy}. None of these
supplies a displacement-matched control. % [EXT]

\paragraph*{Sensitivity and amplification.}
Jacobian and local-linearity methods for transformers are used for steering,
regularization and robustness measures
\cite{skifstad2026local_linearity,martnishn2026dreg_jacobian,zhang2026fisher_robustness,yudin2025attention_lipschitz};
the head split of Section~\ref{sec:adv-reads} is the simplest exact case,
$J{=}W$ at the head input, used to ask what a detector reads. That a small
input perturbation grows into a large internal displacement is also known. The
linear view predicts activation changes that scale with the dimension
\cite{goodfellow2015fgsm}, the Lipschitz constant can grow with depth
\cite{cisse2017parseval}, adversarial perturbations appear as noise in
intermediate feature maps \cite{xie2019feature_denoising}, and a denoiser fails
because the network amplifies the residual perturbation towards the top layers
\cite{liao2018hgd_denoiser}. For linear classifiers, the adversarial direction
moves the output more than random noise of equal norm, by a factor that grows
with the dimension \cite{fawzi2016robustness_adversarial_to_random_noise}. That
work bounds or removes the amplification; this chapter measures it per example
(Section~\ref{sec:adv-amplification}) and asks what a detector that reads the
resulting displacement responds to. % [EXT]

% =============================================================================
\section{The instrument}
\label{sec:adv-instrument}
% Definitions only; no results.

\subsection{Model, data and activations}
\label{sec:adv-models}
The system is a Vision Transformer \emph{classifier}, not a vision--language
model: an ImageNet-1k-pretrained ViT-B/16 backbone
\cite{deng2009imagenet,dosovitskiy2021vit} with a linear classification head
trained on CIFAR-10 \cite{krizhevsky2009cifar} (validation accuracy $0.952$).
% src: config/model/vit_b16.yaml; DESIGN_CHOICES.md §D1 [EST]
CIFAR-10 images are upsampled to $224\times224$ and normalized with ImageNet
statistics.
% src: config/data/cifar10.yaml; DESIGN_CHOICES.md §D1/§D7
ViT-B/16 processes each image as $197$ tokens ($196$ patches and a class token)
through $12$ blocks of width $768$. We probe the MLP outputs of blocks
$\{0,3,6,9,10,11\}$; block~10, \texttt{mlp\_10}, is the probed layer of the
main result. Every activation is mean-pooled over the $197$ tokens, which
limits what the probe can see (Section~\ref{sec:adv-scope}). All runs use
seed~42 unless a five-seed protocol is stated.
% src: config/base.yaml; experiments/2026-06-22_sae-residual-decomp/params.yaml [EST]

\subsection{Notation and the displacement}
\label{sec:adv-dx-axis}
We write $x_i$ for the clean source image $i$ and $\tilde{x}_i$ for a perturbed
version of it, and $a$ for an activation (Section~\ref{sec:adv-sae}). Attack
budgets $\varepsilon$ are in pixel space, with pixel values in $[0,1]$; input
$L_2$ norms are measured in ImageNet-normalized space on the $224\times224$
image, and all perturbations are applied after the upsampling.
% src: src/adv_interp/attacks/fgsm.py:48-50; pgd.py:28-29; data/datamodules.py:33; scripts/calibrate_benign_l2.py:54-62; cli/generate_attacks.py:75-77
For source image $i$ at layer $\ell$, with $a_\ell(\cdot)$ the mean-pooled
activation, the \emph{activation displacement} is
\begin{equation}
  \Delta a_\ell(i) \;=\; a_\ell(\tilde{x}_i) - a_\ell(x_i),
  \label{eq:adv-dx}
\end{equation}
and we measure $\lVert\Delta a\rVert_2$ per example and per layer in every run.
The input norm is a parameter that we set; $\lVert\Delta a\rVert$ is what a
probe at layer $\ell$ can respond to, and the two are related by an
amplification that depends on the family
(Section~\ref{sec:adv-amplification}).

\subsection{Attacks, benign families, and matched controls}
\label{sec:adv-attacks}

We use six attacks, four benign families and one low-displacement control, the
watermark control. All attacks are untargeted and white-box: they use the
gradient of the ViT-B/16 classifier, head included, and they change the input
image only. Table~\ref{tab:adv-attacks} gives the objective and the constraint
of each attack. FGSM at $\varepsilon{=}0.004$ is the reference attack to which
the benign families are matched, except in the PGD replication of
Section~\ref{sec:adv-usual}, where they are matched to PGD at
$\varepsilon{=}0.002$; FGSM and PGD are also swept over
$\varepsilon\in\{0.001,0.002,0.004,0.008,0.012,0.016,0.024,0.032\}$, to test
whether the AUC depends on the budget (Section~\ref{sec:adv-provenance}).
APGD-CE repeats the main result with a stronger attack
(Section~\ref{sec:adv-main}), and the adaptive APGD attacks the classifier and
the probe together (Section~\ref{sec:adv-attackers}). The probe is a logistic
regression on standardized, mean-pooled activations
(Section~\ref{sec:adv-protocol}). It is differentiable and smooth, and it has
no random part. Thus, gradient masking is not probable, and methods for masked
gradients such as BPDA and EOT
\cite{athalye2018obfuscated_gradients,tramer2020adaptive_attacks} are not
necessary. We measure the realized input norm and the attack success rate
(ASR) in every run; the ASR is the fraction of the images classified correctly
when clean ($1536$ of $1600$ in the main run) whose label the attack changes:
FGSM has $L_2=6.85$ and ASR $0.762$ at $\varepsilon{=}0.004$, PGD has ASR
$0.922$ at $\varepsilon{=}0.002$, and CW-$L_2$ has ASR $1.000$.
% src: 2026-06-29_h2-v2-eps0004/outputs/h2_crossval/metrics.json (l2_report); 2026-07-02_pgd-h2-eps0002/outputs/h2_crossval/metrics.json; 2026-07-21_cw-attack/outputs/full/asr_cw.json [EST]

\begin{table}[ht]
  \centering
  \caption{Attacks. The input norm is either set by the budget $\varepsilon$ or
  measured after the attack. The adaptive APGD and the minimal-displacement
  attack are defined in Section~\ref{sec:adv-attackers}.}
  \label{tab:adv-attacks}
  \small
  \begin{tabular}{@{}>{\raggedright\arraybackslash}p{0.18\linewidth}>{\raggedright\arraybackslash}p{0.16\linewidth}>{\raggedright\arraybackslash}p{0.24\linewidth}>{\raggedright\arraybackslash}p{0.15\linewidth}>{\raggedright\arraybackslash}p{0.11\linewidth}@{}}
    \toprule
    Attack & Steps & Objective & Constraint & Input norm \\
    \midrule
    FGSM \cite{goodfellow2015fgsm} & $1$ & Increase the classification loss & $L_\infty\le\varepsilon$ & Set by $\varepsilon$ \\
    PGD \cite{madry2018pgd} & $20$, step $2.5\varepsilon/20$, random start & Increase the classification loss & $L_\infty\le\varepsilon$ & Set by $\varepsilon$ \\
    APGD-CE \cite{croce2020apgd_autoattack} & $100$, $1$ restart & Increase the classification loss & $L_\infty\le\varepsilon$ & Set by $\varepsilon$ \\
    Adaptive APGD & $100$ per $\lambda$, $1$ restart & Wrong label and probe score below $\tau$ & $L_\infty\le\varepsilon$ & Set by $\varepsilon$ \\
    CW-$L_2$ \cite{carlini2017cw_attack} & $50$, fixed $c{=}1$ & Decrease the input $L_2$ plus a margin term & Margin penalty ($\kappa{=}0$) & Measured \\
    Minimal-displacement & Iterative, binary search on $c$ & Decrease $\lVert\Delta a\rVert$ at \texttt{mlp\_10} & Label must change & Measured \\
    \bottomrule
  \end{tabular}
\end{table}

\paragraph*{Benign families and two matched controls.}
To separate ``adversarial'' from ``perturbed'', we use four benign perturbation
families: Gaussian noise, Gaussian blur, JPEG compression, and random pixel
dropout. Each family has one parameter, which we calibrate so that the family
matches FGSM in one of two matched controls (Table~\ref{tab:adv-controls}).
The displacement-matched control is exploratory
(Section~\ref{sec:adv-provenance}). Non-matched variants of the benign
families at $L_2\approx17$--$21$ reach AUC $0.853$--$0.938$ on $z$; they are
never used in a decision rule.
% src: 2026-06-29_h2-v2-eps0004 metrics.json; 2026-07-22_h2-matched-dx/outputs/h2_matched_dx_summary.json (realized_mean_dx_mlp10); calib/dx_calibration.json [EST — E1]

\begin{table}[ht]
  \centering
  \caption{The two matched controls. Both match the benign families to FGSM at
  $\varepsilon{=}0.004$; the PGD replication of the input-matched control
  (Section~\ref{sec:adv-usual}) matches them to PGD at $\varepsilon{=}0.002$.}
  \label{tab:adv-controls}
  \small
  \begin{tabular}{@{}>{\raggedright\arraybackslash}p{0.17\linewidth}>{\raggedright\arraybackslash}p{0.37\linewidth}>{\raggedright\arraybackslash}p{0.37\linewidth}@{}}
    \toprule
     & Input-matched control & Displacement-matched control \\
    \midrule
    Matched quantity & Per-image input $L_2$, within $\pm20\%$ & Mean $\lVert\Delta a\rVert$ at \texttt{mlp\_10} (target $1.994$, the FGSM mean) \\
    Method & Family parameter set from the $L_2$ target & Model-in-the-loop search on the family parameter \\
    Verification & Ratios $0.999$--$1.008$ & Realized $1.78$--$2.06$ (calibration), $1.78$--$2.11$ (evaluation) \\
    Granularity & Per image & Mean over images \\
    Layers & Not layer-specific (input space) & Calibrated at \texttt{mlp\_10}; other layers inherit it \\
    Status & Pre-specified & Exploratory, expectation written before the run \\
    \bottomrule
  \end{tabular}
\end{table}

\paragraph*{The watermark control.}
A faint text watermark (steganography, opacity $0.03$) has an input $L_2$ of
$2.24$ ($0.33\times$ FGSM at $\varepsilon{=}0.004$) and the smallest
displacement of any family ($\lVert\Delta a\rVert=0.09$ at \texttt{mlp\_10},
under $1\%$ of the clean activation norm). It changes almost no label (ASR
$\approx0.0007$), so it is not adversarial in the functional sense.
% src: experiments/2026-06-29_h2-v2-eps0004/outputs/h2_crossval/metrics.json (asr block) [EST]
It anchors the low-displacement
column of the design matrix (Table~\ref{tab:adv-design}) by measurement, not
by assumption.
% src: 2026-06-29_h2-v2-eps0004 metrics.json (l2_report); 2026-07-23_jpeg-matched-l2-dx/outputs/displacement_summary.json [EST]

\paragraph*{The design matrix.}
Table~\ref{tab:adv-design} crosses the functional effect (does the
perturbation change the label?) with the activation displacement. FGSM, PGD
and CW-$L_2$ are in the adversarial, high-displacement cell; the benign
families are in the non-adversarial row, at low displacement in the
input-matched control and at high displacement in the displacement-matched
control; and the watermark control is at low displacement. No standard attack
occupies the adversarial, low-displacement cell;
Section~\ref{sec:adv-attackers} occupies it by construction.
% NOTE: the non-adversarial row is labelled "not optimized against the
% classifier" because the rate at which the displacement-matched benign
% families change the label was not measured (paper action point A2).

\begin{table}[ht]
  \centering
  \caption{Design matrix. In the two rows, a probe trained on a family at high
  displacement detects it. The adversarial cell at low displacement is the
  exception: the probe detects these examples, and the cause is an open
  question (Section~\ref{sec:adv-attackers}).}
  \label{tab:adv-design}
  \small
  \renewcommand{\arraystretch}{1.25}
  \begin{tabular}{@{}>{\raggedright\arraybackslash}p{0.17\linewidth}>{\raggedright\arraybackslash}p{0.37\linewidth}>{\raggedright\arraybackslash}p{0.37\linewidth}@{}}
    \toprule
     & \textbf{Low $\lVert\Delta a\rVert$} & \textbf{High $\lVert\Delta a\rVert$} \\
    \midrule
    \textbf{Adversarial} \newline {\footnotesize label changes}
     & \textbf{Empty for standard attacks} \newline
       Minimal-displacement attack \newline
       $\lVert\Delta a\rVert = 0.298$, ASR $1.000$ \newline
       Input $L_\infty = 0.137$, $L_2$ $1.22\times$ FGSM \newline
       Probe AUC $0.999$ (open question)
     & \textbf{FGSM, PGD, CW-$L_2$} \newline
       Optimize misclassification only \newline
       PGD grid: $\lVert\Delta a\rVert$ $1.24$ to $3.11$ \newline
       FGSM $2.00$, CW-$L_2$ $1.81$ \newline
       Detected; AUC rises with $\varepsilon$ \\
    \midrule
    \textbf{Not optimized} \newline {\footnotesize against the classifier}
     & \textbf{Watermark, benign at matched $L_2$} \newline
       Watermark: $\lVert\Delta a\rVert = 0.09$, AUC $0.542$ \newline
       Benign at matched input $L_2$: \newline
       $\lVert\Delta a\rVert$ $0.19$ to $0.59$ \newline
       Less detectable than FGSM
     & \textbf{Benign at matched $\lVert\Delta a\rVert$} \newline
       Noise, blur, JPEG, dropout \newline
       $\lVert\Delta a\rVert$ $1.78$ to $2.11$ (target $1.994$) \newline
       At least as detectable as FGSM \newline
       $\dauc \le 0$ at \texttt{mlp\_10}, five seeds \\
    \bottomrule
  \end{tabular}
\end{table}

\subsection{Sparse autoencoders and the residual decomposition}
\label{sec:adv-sae}
The SAE decomposes the activation as
\begin{equation}
  a \;=\; \hat{a} + r,
  \qquad \hat{a} = \mathrm{dec}(z), \qquad z = \mathrm{ReLU}(\mathrm{enc}(a)),
  \label{eq:adv-decomp}
\end{equation}
where $a$ is the mean-pooled raw activation, $z$ the sparse latents, $\hat{a}$
the reconstruction, and $r=a-\hat{a}$ the residual. We use a vanilla $L_1$ SAE
(input dimension $768$, $1024$ latents, $L_1$ penalty $0.1$; Adam, learning
rate $10^{-3}$, batch $256$, ten epochs). % src: config/sae/default.yaml [EST]
It is trained on \emph{clean activations only}, the $1600$ mean-pooled clean
activations of the run, so $r$ is the part of the activation that the clean
dictionary cannot explain. These clean activations are also the clean class of
every detection task, and no clean images are held out from the SAE; this can
raise the AUCs on $z$, $\hat{a}$ and $r$, and it does not affect the AUCs on
$a$. We evaluate detection on the four vector representations $a$, $z$,
$\hat{a}$, $r$ and on three scalar norms ($\lVert r\rVert_2$,
$\lVert\hat{a}\rVert_2$, $\lVert z\rVert_1$).
% src: experiments/2026-06-22_sae-residual-decomp/params.yaml [EST]
% NOTE: the July version said that \hat{a} "projects onto the clean-data
% manifold" and that r is "the off-manifold part". The residual of an SAE
% trained on clean data is not a projection onto a manifold; the wording was
% dropped here (the paper still carries it in its SAE appendix: action point B12).

\paragraph*{SAE health gates.}
Every comparison on $z$, $\hat{a}$ or $r$ is gated on the health of the SAE:
fraction of variance explained (FVE) $\geq 0.9$ and dead-feature fraction
$\leq 30\%$; we also report $L_0$, which has no threshold. % [PRE — health gate]
The SAE at \texttt{mlp\_10} is healthy (FVE $0.938$, $L_0\approx135$, dead
fraction $3$--$5\%$); away from it, health varies by run, and in practice only
\texttt{mlp\_10} and \texttt{mlp\_11} pass (early layers reach FVE
$0.078$--$0.352$ in the six-layer decomposition).
% src: magnitude_sweep_metrics.csv; h2_matched_dx_summary.json; 2026-07-21_decomp-6layer [EST]
% Internal note carried over from the paper: FVE at mlp_0 differs between runs
% (0.75 in the 2026-06-22 decomposition vs 0.08 in the 2026-07-21 six-layer run);
% unexplained fit difference, flagged rather than resolved.
Claims about $z$, $\hat{a}$ and $r$ are therefore restricted to these two
layers; claims about the raw activation $a$ are exempt.

\subsection{Detection protocol and statistical controls}
\label{sec:adv-protocol}
Every detection AUC comes from one routine. Each binary task compares clean
images with \emph{one} family, never a pool, at $N{=}3200$ ($1600$ clean,
$1600$ perturbed). We use grouped $K$-fold cross-validation ($K{=}5$) by
source image: all versions of image~$i$ (clean, attacked, benign and
watermark) are in the same fold, so the probe is never tested on a near-copy
of a training image. Features are standardized on the training fold only, and
the probe is a logistic regression (LBFGS, at most $1000$ iterations).
% src: src/adv_interp/detection/evaluation.py; experiments/2026-06-29_h2-v2-eps0004/outputs/h2_crossval/metrics.json [EST]

\paragraph*{What a linear probe can read.}
Because the probe is linear, the change of its score has a simple form. Let
$w$ be the probe weights, $\sigma_k$ the training-fold standard deviation of
feature $k$, and $\tilde{w}_k = w_k/\sigma_k$. The logit $s$ of a perturbed
image then differs from the logit of its clean source by
\begin{equation}
  s(\tilde{x}_i) - s(x_i) \;=\; \tilde{w}^{\top}\Delta a
  \;=\; \lVert\tilde{w}\rVert\,\lVert\Delta a\rVert\cos\theta,
  \label{eq:adv-probe}
\end{equation}
where $\theta$ is the angle between $\Delta a$ and $\tilde{w}$. The probe
therefore reads the component of $\Delta a$ along one learned direction:
$\lVert\Delta a\rVert$ sets the scale of this component, and $\cos\theta$ sets
its sign and the fraction of the displacement that the probe sees. Two
families with equal $\lVert\Delta a\rVert$ can thus get different scores. The
same form holds for each representation of Equation~\eqref{eq:adv-decomp},
with $\Delta a$ replaced by the change of that representation.
Equation~\eqref{eq:adv-probe} describes the shift of one example; the AUC also
depends on how these shifts are spread over examples.

\paragraph*{The detection advantage.}
Throughout, specificity is summarized by the signed quantity
\begin{equation}
  \dauc \;=\; \mathrm{AUC}(\text{clean vs.\ attack})
  \;-\; \max_{b}\;\mathrm{AUC}(\text{clean vs.\ benign}_{b}),
  \label{eq:adv-dauc}
\end{equation}
with the maximum taken over the magnitude-matched benign families $b$. It is
positive when the attack is more detectable than every benign perturbation of
equal magnitude, and negative when some benign family is more detectable than
the attack; nothing constrains its sign. When $\dauc>0$, we say that the
attack has a \emph{detection advantage}. Each AUC in $\dauc$ comes from a
probe trained for its own clean-vs-family task, and
\pendingpaper{A3}{we did not train an attack-vs-benign probe.}

\paragraph*{Statistical controls.}
The five folds share one seed, one SAE and one attack draw, so the standard
deviation over folds understates the variance. A measurement bug that gave a
false flat sweep (Section~\ref{sec:adv-bugs}) also showed that the pipeline
must be seen to return chance when it should. Four controls calibrate what
the AUCs mean. A clean-vs-clean null gives AUC $0.500\pm0.018$ ($10$ repeats).
A within-group label-permutation null gives $0.500\pm0.013$, which places the
observed FGSM AUC of $0.932$ at $\varepsilon{=}0.004$ at $p=0.0099$, the
smallest value that $100$ permutations can give.
% src: experiments/2026-07-21_stat-package/outputs/stat_package.json (b, c blocks) [EST — E4b/E4c]
The headline $\dauc$ values also carry a paired bootstrap over source images
($95\%$ CI) and a five-seed replication, with the SAE retrained and the
perturbations regenerated for each seed
(Sections~\ref{sec:adv-usual} and~\ref{sec:adv-main}).
% src: same file (a block); experiments/2026-07-02_gap-robustness/outputs/gap_robustness_summary.csv [EST — E4a/G5.3]
We read specificity in the range before saturation
($\varepsilon\approx0.002$--$0.004$): FGSM reaches AUC $1.000$ at
$\varepsilon\geq0.012$, and at the ceiling no comparison can separate the two
readings of the detector. % src: magnitude_sweep_metrics.csv

% =============================================================================
\section{Step 1: the usual control shows a detection advantage}
\label{sec:adv-usual}
\emph{Objection: ``The input-norm match is the standard control, and it shows
a detection advantage.''}

The standard control does report a detection advantage, and a large one. In
the input-matched control at $\varepsilon{=}0.004$, FGSM is the most
detectable family on $z$ and on $a$ (Table~\ref{tab:adv-main}): $\dauc =
0.199$ on $z$ at \texttt{mlp\_10}, with a paired-bootstrap $95\%$ CI of
$[0.188, 0.212]$ and $0.204\pm0.014$ over five seeds. PGD at
$\varepsilon{=}0.002$ gives the same picture. For this replication only, the
benign families are matched to PGD in input $L_2$ ($3.06$); JPEG is absent,
because its $L_2$ cannot reach the $\pm20\%$ band. $\dauc = 0.224$, with every
benign family below $0.61$. APGD-CE at $\varepsilon{=}0.004$ also gives the
same picture, with $\dauc = 0.271\pm0.010$ over five seeds
(Section~\ref{sec:adv-main}).
% src: 2026-06-29_h2-v2-eps0004 and 2026-07-02_pgd-h2-eps0002 metrics.json; stat_package.json; gap_robustness_summary.csv [EST]
Without the displacement-matched control, this is the number that one would
read as adversarial specificity, and it is the strongest evidence for $H_{M}$
in the study.

The input match is verified in the run, not assumed from the parameters of the
families (Table~\ref{tab:adv-input-match}): the realized per-image input $L_2$
of each benign family is within $0.999$--$1.008$ of that of FGSM, and within
$0.939$--$1.001$ of that of PGD.

\begin{table}[ht]
  \centering
  \caption{Input-matched control, probe on $z$ at \texttt{mlp\_10}: detection
  AUC of each family and the ratio of its measured input $L_2$ to that of the
  attack. JPEG is absent from the PGD block because its $L_2$ cannot reach the
  $\pm20\%$ band. The interval is a paired-bootstrap $95\%$ CI over source
  images.}
  \label{tab:adv-input-match}
  \small
  \begin{tabular}{@{}lcc@{}}
    \toprule
    Family & AUC on $z$ & $L_2$ ratio \\
    \midrule
    \multicolumn{3}{@{}l}{\emph{FGSM, $\varepsilon{=}0.004$}}\\
    FGSM                  & $0.932$ & --- \\
    Gaussian blur         & $0.733$ & $1.000$ \\
    Dropout               & $0.617$ & $1.005$ \\
    JPEG                  & $0.604$ & $1.008$ \\
    Gaussian noise        & $0.574$ & $0.999$ \\
    $\dauc$               & $0.199$ $[0.188, 0.212]$ & \\
    \midrule
    \multicolumn{3}{@{}l}{\emph{PGD, $\varepsilon{=}0.002$}}\\
    PGD                   & $0.827$ & --- \\
    Gaussian blur         & $0.603$ & $1.000$ \\
    Dropout               & $0.524$ & $0.939$ \\
    Gaussian noise        & $0.512$ & $1.001$ \\
    $\dauc$               & $0.224$ & \\
    \bottomrule
  \end{tabular}
  % src: experiments/2026-06-29_h2-v2-eps0004/outputs/h2_crossval/metrics.json;
  %      experiments/2026-07-02_pgd-h2-eps0002/outputs/h2_crossval/metrics.json;
  %      experiments/2026-07-21_stat-package/outputs/stat_package.json
  % The July version printed "+/-" values next to each AUC, carried from the
  % paper without a stated meaning; they are omitted until their meaning is
  % confirmed (fold standard deviations?).
\end{table}

\begin{table}[ht]
  \centering
  \caption{Main result at \texttt{mlp\_10}, FGSM at $\varepsilon{=}0.004$. Mean
  displacement and detection AUC of the clean-vs-family probe on the raw
  activation $a$ and on the SAE latents $z$, in the two matched controls. Bold:
  the most detectable benign family, which sets $\dauc$
  (Equation~\eqref{eq:adv-dauc}); in the displacement-matched control on $a$,
  Gaussian noise ($1.0000$) is above JPEG ($0.9999$). Each control is a
  separate run, with the FGSM examples regenerated and the probe retrained, so
  FGSM differs between the two controls by at most $0.001$. Five-seed values
  are mean $\pm$ standard deviation, with the SAE retrained and the
  perturbations regenerated for each seed.}
  \label{tab:adv-main}
  \footnotesize
  \setlength{\tabcolsep}{3pt}
  \begin{tabular}{@{}lcccccc@{}}
    \toprule
     & \multicolumn{3}{c}{Input-matched control} & \multicolumn{3}{c}{Displacement-matched control} \\
    \cmidrule(lr){2-4}\cmidrule(lr){5-7}
    Family & $\lVert\Delta a\rVert$ & AUC on $a$ & AUC on $z$ & $\lVert\Delta a\rVert$ & AUC on $a$ & AUC on $z$ \\
    \midrule
    FGSM & $2.00$ & $0.989$ & $0.932$ & $2.00$ & $0.989$ & $0.933$ \\
    Gaussian noise & $0.19$ & $0.673$ & $0.574$ & $2.01$ & $\mathbf{1.000}$ & $0.994$ \\
    Gaussian blur & $0.59$ & $\mathbf{0.923}$ & $\mathbf{0.733}$ & $1.78$ & $0.998$ & $0.954$ \\
    JPEG & $0.28$ & $0.721$ & $0.604$ & $2.06$ & $1.000$ & $\mathbf{0.995}$ \\
    Dropout & $0.21$ & $0.705$ & $0.617$ & $2.11$ & $0.999$ & $0.989$ \\
    \midrule
    $\dauc$, seed 42 & & $+0.066$ & $+0.199$ & & $-0.011$ & $-0.062$ \\
    $\dauc$, five seeds & & $+0.065\pm0.006$ & $+0.204\pm0.014$ & & $-0.013\pm0.002$ & $-0.063\pm0.003$ \\
    \bottomrule
  \end{tabular}
  % src: figures/tab_main.tex in the paper repository (generated by figures/src/tab_main.py)
\end{table}

% =============================================================================
\section{Step 2: the usual control is not matched in activation displacement}
\label{sec:adv-amplification}
\emph{Objection: ``The benign perturbations have the same input norm as the
attack, so they have the same magnitude.''}

The input norm does not set the displacement
(Figure~\ref{fig:adv-amplification}). At matched input $L_2$
($\varepsilon{=}0.004$), FGSM arrives at \texttt{mlp\_10} with
$\lVert\Delta a\rVert=2.00$, and the four benign families arrive with
$0.19$--$0.59$ (Table~\ref{tab:adv-main}).
% src: experiments/2026-07-21_attack-lib-smoke/outputs/smoke_results.json; experiments/2026-07-23_jpeg-matched-l2-dx/outputs/displacement_summary.json [EST]
The clean activation norm at that layer is $10.0$, so the attack moves the
activation by $20\%$ of its length and the benign families move it by
$2$--$6\%$. CW-$L_2$ shows the same relation from the other side. Because it
minimizes the input norm, it was the planned candidate for the adversarial,
low-displacement cell of Table~\ref{tab:adv-design}.
% src: EXPERIMENT_PLAN_ATTACKS.md §0
It changes
every label (ASR $1.000$) with an input $L_2$ of $4.97$, $0.72\times$ that of
FGSM, and it still arrives with a mean $\lVert\Delta a\rVert$ of $1.81$
(standard deviation over images $1.39$) and is detected at AUC $0.865$.
% src: experiments/2026-07-21_cw-attack/outputs/full/cw_input_l2.json; asr_cw.json; displacement_summary.json; eval/metrics.json [EST — E7]
PGD follows the same relation along its budget grid, from
$\lVert\Delta a\rVert=1.24$ at $\varepsilon{=}0.001$ to $3.11$ at $0.032$; PGD
at $\varepsilon{=}0.002$ and FGSM at $0.004$ have the same displacement
($2.01$ and $2.00$).
% src: experiments/2026-07-21_pgd-displacement/outputs/pgd_displacement_summary.csv [EST — E6]

We call the ratio of displacement to input norm the \emph{amplification}. It
is much larger for perturbations made by adversarial optimization, which
follows the gradient of the network, than for the benign families. Every
attack that optimizes misclassification alone thus lands in the
high-displacement column of Table~\ref{tab:adv-design}. This is where
adversarial optimization enters the account: as a cause of the displacement,
before anything that the probe reads.

\begin{figure}[ht]
  \centering
  \includegraphics[width=0.7\linewidth]{figures/interp_adv/fig_amplification.pdf}
  \caption{Input norm and displacement at \texttt{mlp\_10}. At the same input
  $L_2$, FGSM arrives with $3$--$10\times$ the displacement of the benign
  families; CW-$L_2$ changes every label with a smaller input norm and still
  arrives with an attack-scale displacement; PGD arrives with $2.01$ at
  $\varepsilon{=}0.002$ and $2.50$ at $\varepsilon{=}0.004$; the watermark
  control almost does not move the activation.}
  % src: figures/src/fig_amplification.py (paper repository)
  \label{fig:adv-amplification}
\end{figure}

% =============================================================================
\section{Step 3: matched displacement removes the detection advantage}
\label{sec:adv-main}
\emph{Objection: ``The displacement match is the wrong control.''}

The probe reads $\Delta a$ at its layer (Equation~\eqref{eq:adv-probe}) and
does not see the input norm, so the quantity to match is
$\lVert\Delta a\rVert$. The displacement-matched control is the decisive one.
We calibrate the four benign families to FGSM's displacement at
\texttt{mlp\_10} (Section~\ref{sec:adv-attacks}) and train the clean-vs-family
probes again at every layer. Table~\ref{tab:adv-main} gives the AUCs at
\texttt{mlp\_10} in the two controls. In the displacement-matched control,
every benign family is at least as detectable as FGSM, on $z$ and on $a$, so
$\dauc$ is negative. The matched benign families are strong corruptions (for
example, blur with $\sigma=8.00$ and JPEG quality $17$;
Table~\ref{tab:adv-params}), and most of their AUCs are near the ceiling. At
\texttt{mlp\_11}, $\dauc=-0.078$ on $z$ (dropout).
% src: 2026-07-22_h2-matched-dx/outputs/h2_matched_dx_summary.json; decomposition_metrics.csv; figures/src/tab_main.py [EST — E1]

A five-seed replication retrains the SAE and regenerates the perturbations for
each seed, and it verifies the displacement match in all $25$ family--seed
cells (the four benign families and FGSM, in five seeds). On $z$, the mean and
the $+2\sigma$ upper bound of $\dauc$ stay negative at \texttt{mlp\_10}
(Table~\ref{tab:adv-main}) and at \texttt{mlp\_11} ($-0.076\pm0.005$). On $a$,
the mean $\dauc$ is zero or negative at all six layers, from $-0.032\pm0.001$
at \texttt{mlp\_9} to $0.000\pm0.000$ at \texttt{mlp\_3}.
% src: experiments/2026-07-23_h2-matched-dx-multiseed/outputs/multiseed_summary.csv; multiseed_verdict.json; README.md (caveat) [EST — replicated 2026-07-23]

The detection advantage of $0.199$ in the input-matched control therefore
needs no adversarial meaning to explain it: it is the amplification of
Section~\ref{sec:adv-amplification}, read by a control that is matched in the
wrong quantity.

\subsection{The match at each layer}
\label{sec:adv-dx-layers}
The displacement-matched control is calibrated at \texttt{mlp\_10} only, so
``$\dauc\le0$ at every layer'' would overstate the result.
Table~\ref{tab:adv-dx-layers} gives the realized mean displacement of each
family at each probed layer, and Table~\ref{tab:adv-params} gives the
calibrated parameter of each benign family in the two matched controls. At
\texttt{mlp\_9} and \texttt{mlp\_11} every benign family has a smaller mean
displacement than FGSM ($0.64$--$0.97\times$), and it is still at least as
detectable; these layers agree with the result at \texttt{mlp\_10}. At
\texttt{mlp\_0} and \texttt{mlp\_3} every benign family has a larger mean
displacement than FGSM ($1.1$--$10.5\times$), so $\dauc\le0$ at these layers
is not evidence about equal displacement. In the five-seed replication, one
flag applies: the dead-feature fraction exceeds $0.30$ at \texttt{mlp\_11} in
two seeds, and excluding these seeds does not change the verdict.
% src: experiments/2026-07-22_h2-matched-dx/outputs/displacement.csv; h2_matched_dx_summary.json (per_layer gap_x); calib/dx_calibration.json; config/attack/benign_*_l2_0004.yaml

\begin{table}[ht]
  \centering
  \caption{Displacement-matched control, seed 42: mean $\lVert\Delta a\rVert$
  of each family at each probed layer ($1600$ examples per cell), with the
  ratio to FGSM in parentheses, and $\dauc$ on the raw activation $a$. The
  control is calibrated at \texttt{mlp\_10}.}
  \label{tab:adv-dx-layers}
  \footnotesize
  \setlength{\tabcolsep}{3pt}
  \begin{tabular}{@{}lcccccc@{}}
    \toprule
    Layer & FGSM & Noise & Blur & JPEG & Dropout & $\dauc$ on $a$ \\
    \midrule
    \texttt{mlp\_0} & $0.09$ & $0.91$ ($10.5\times$) & $0.21$ ($2.4\times$) & $0.33$ ($3.9\times$) & $0.86$ ($10.0\times$) & $-0.001$ \\
    \texttt{mlp\_3} & $0.97$ & $2.47$ ($2.5\times$) & $1.07$ ($1.1\times$) & $1.92$ ($2.0\times$) & $2.42$ ($2.5\times$) & $\phantom{-}0.000$ \\
    \texttt{mlp\_6} & $1.31$ & $1.42$ ($1.1\times$) & $1.15$ ($0.9\times$) & $2.06$ ($1.6\times$) & $1.38$ ($1.1\times$) & $-0.007$ \\
    \texttt{mlp\_9} & $1.47$ & $1.02$ ($0.70\times$) & $0.94$ ($0.64\times$) & $1.38$ ($0.94\times$) & $1.00$ ($0.68\times$) & $-0.032$ \\
    \texttt{mlp\_10} & $2.00$ & $2.01$ ($1.00\times$) & $1.78$ ($0.89\times$) & $2.06$ ($1.03\times$) & $2.11$ ($1.05\times$) & $-0.011$ \\
    \texttt{mlp\_11} & $1.08$ & $0.93$ ($0.87\times$) & $0.91$ ($0.84\times$) & $1.05$ ($0.97\times$) & $0.91$ ($0.84\times$) & $-0.011$ \\
    \bottomrule
  \end{tabular}
\end{table}

\begin{table}[ht]
  \centering
  \caption{Calibrated parameter of each benign family in the two matched
  controls (FGSM at $\varepsilon{=}0.004$, seed 42).}
  \label{tab:adv-params}
  \small
  \begin{tabular}{@{}llcc@{}}
    \toprule
    Family & Parameter & Input-matched control & Displacement-matched control \\
    \midrule
    Gaussian noise & Standard deviation $\sigma$ & $0.0040$ & $0.0266$ \\
    Gaussian blur & Standard deviation $\sigma$ & $1.17$ & $8.00$ \\
    JPEG & Quality & $93$ & $17$ \\
    Dropout & Pixel probability $p$ & $6.2\times10^{-5}$ & $3.1\times10^{-3}$ \\
    \bottomrule
  \end{tabular}
\end{table}

\subsection{A stronger attack}
\label{sec:adv-apgd-ce}
FGSM is a weak attack, so we repeat the two matched controls with APGD-CE
\cite{croce2020apgd_autoattack} at the same budget ($\varepsilon{=}0.004$; ASR
$1.000$; five seeds). The test was pre-specified, with the decision rule of
the five-seed replication: the attack has a detection advantage if the mean
minus $2\sigma$ of $\dauc$ is above $0.05$, and none if the mean plus
$2\sigma$ is below $0.05$. APGD-CE arrives at \texttt{mlp\_10} with a mean
$\lVert\Delta a\rVert$ of $2.70$ on the calibration images, $1.35\times$ that
of FGSM, with a smaller input $L_2$ ($6.03$ against $6.85$ in seed $42$). In
the input-matched control we use the benign families matched to FGSM (their
input $L_2$ is $1.12$--$1.15\times$ that of APGD-CE); in the
displacement-matched control we calibrate them again to
$\lVert\Delta a\rVert = 2.70$ (realized ratios $0.96$--$1.05$).
% src: experiments/2026-10-19_apgd-p7/outputs/pilot/seed42/pilot_metrics.json; calib/dx_calibration.json; metrics.json (matching_per_seed); JOURNAL_V3.md §8b [PRE — P7]

In the input-matched control the attack again has a detection advantage:
$\dauc = 0.077\pm0.005$ on $a$ and $0.271\pm0.010$ on $z$ at \texttt{mlp\_10}.
In the displacement-matched control it disappears ($\dauc = -0.001\pm0.000$ on
$a$ and $-0.003\pm0.001$ on $z$), but at the ceiling: the most detectable
benign family (Gaussian noise) is at AUC $1.000$ on $a$ and $0.999$ on $z$, so
$\dauc$ cannot be much above zero.
% src: experiments/2026-10-19_apgd-p7/outputs/metrics.json (gated_cells); p7_cells.csv [PRE — P7 branch B]
\pendingpaper{A1}{The evidence that the advantage disappears at equal
displacement therefore comes from FGSM on $z$, where FGSM ($0.933$) and blur
($0.954$) are below the ceiling; a replication with APGD-CE below the ceiling
is not available.}

\subsection{The result is not a property of the SAE, and pooling bounds it}
\label{sec:adv-scope}
\emph{Objection: ``The result is a property of the SAE or of the pooling.''}

\paragraph*{The SAE is not needed for detection.}
At $\varepsilon{=}0.004$, the raw activation $a$ detects FGSM at $0.989$, and
the latents $z$ at $0.932$; for detection, this SAE loses signal.
% src: experiments/2026-06-30_decomp-eps0004/outputs/decomposition_metrics.csv [EST — L2]
In the displacement-matched control, the sign of $\dauc$ is the same on $a$
and on $z$ (Table~\ref{tab:adv-main}); on $a$, the comparison is at the
ceiling. In the input-matched control, $\dauc$ is larger on the SAE-explained
component ($0.199$ on $z$ and $0.178$ on $\hat{a}$) than on $r$ ($0.079$) and
on $a$ ($0.066$; Table~\ref{tab:adv-representations}). This does not show that
the SAE amplifies specificity. The probes on $a$ and on $r$ detect FGSM near
the ceiling, where a difference of AUCs is compressed, and a pre-specified
anti-condition forbids the claim: at $\varepsilon{=}0.004$ the difference on
the logit scale on $a$ is compatible with that on $z$, and only at
$\varepsilon{=}0.002$ is $z$ distinct from $a$.
% src: gap_robustness_summary.csv (logit columns); DESIGN_CHOICES.md D12 (G5.3, frozen 2026-07-02) [PRE — anti-condition]
% NOTE: action point B5 of the paper ("the SAE inflates the apparent
% specificity", 0.199 on z vs 0.066 on a) conflicts with this anti-condition.
The input-matched advantage also depends on depth: it is negative at \texttt{mlp\_0} and rises
to $0.204\pm0.014$ at \texttt{mlp\_10} over five seeds, while detection of the
displacement is already present at the first block.
% src: experiments/2026-07-02_gap-robustness/outputs/gap_robustness_summary.csv [EST]

\begin{table}[ht]
  \centering
  \caption{Detection AUC per representation at \texttt{mlp\_10}, input-matched
  control, $\varepsilon{=}0.004$, one seed (FVE $0.938$). The vector
  representations detect the displacement; the scalar norms are near chance.
  $\dauc$ is FGSM minus the most detectable benign family.}
  \label{tab:adv-representations}
  \small
  \begin{tabular}{@{}lccccccc@{}}
    \toprule
    Family & $a$ & $z$ & $\hat{a}$ & $r$ & $\lVert r\rVert$ & $\lVert\hat{a}\rVert$ & $\lVert z\rVert_1$ \\
    \midrule
    FGSM & $0.989$ & $0.932$ & $0.963$ & $0.990$ & $0.594$ & $0.570$ & $0.578$ \\
    Noise & $0.673$ & $0.574$ & $0.608$ & $0.660$ & $0.505$ & $0.512$ & $0.510$ \\
    Blur & $0.923$ & $0.733$ & $0.785$ & $0.911$ & $0.540$ & $0.557$ & $0.552$ \\
    JPEG & $0.721$ & $0.604$ & $0.653$ & $0.698$ & $0.501$ & $0.515$ & $0.514$ \\
    Dropout & $0.705$ & $0.617$ & $0.659$ & $0.687$ & $0.516$ & $0.507$ & $0.506$ \\
    \midrule
    $\dauc$ & $+0.066$ & $+0.199$ & $+0.178$ & $+0.079$ & $+0.054$ & $+0.013$ & $+0.026$ \\
    \bottomrule
  \end{tabular}
  % src: figures/tab_representations.tex in the paper repository; experiments/2026-06-30_decomp-eps0004/outputs/decomposition_metrics.csv
\end{table}

\paragraph*{Pooling hides localized signal.}
Mean-pooling over $197$ tokens can remove a localized perturbation, and the
pooled probe does not detect the watermark control ($0.542$ on $a$, $0.519$ on
$z$ at \texttt{mlp\_10}). A token-by-token probe finds the watermark tokens:
the best single token reaches AUC $0.667\pm0.010$ (standard deviation over the
five folds; Figure~\ref{fig:adv-token}), and tokens $1$ to $6$ span
$0.601$--$0.667$, against $0.542$ for the mean-pooled probe, $0.517$ for the
class token alone and $0.512$ for a max-pooled probe.
% src: experiments/2026-07-21_stego-token/outputs/token_probe_metrics.json [EST — E9]
The watermark is therefore invisible to the pooled probe because its pooled
displacement is small ($\lVert\Delta a\rVert=0.09$), not because the
activations carry no trace of it. This limits every null result of the chapter
to pooled probes (Section~\ref{sec:adv-closing}).

\begin{figure}[ht]
  \centering
  \includegraphics[width=0.5\linewidth]{figures/interp_adv/fig_token.pdf}
  \caption{Per-token detection AUC of the watermark control at
  \texttt{mlp\_10} (probe on $a$, one token at a time; the $14\times14$ grid of
  image patches). The watermark tokens at the top of the image are detected
  (best token $0.667$), while the mean-pooled probe reaches $0.542$ and the
  class token $0.517$.}
  % src: figures/src/appendix_items.py; 2026-07-21_stego-token/outputs/token_probe_metrics.json
  \label{fig:adv-token}
\end{figure}

% =============================================================================
\section{Step 4: the probe reads a direction, not the change of the label}
\label{sec:adv-reads}
\emph{Objection: ``Matching $\lVert\Delta a\rVert$ removes exactly the
adversarial signal, because adversarial examples are defined by a large
displacement.''}

If this objection were correct, the probe would respond to the property that
makes an example adversarial: the change of the label. Two experiments show
that it does not.

\subsection{Detection does not depend on attack success}
\label{sec:adv-success}
We split the FGSM examples by whether the attack changed the label. At every
budget below saturation, successful and failed attacks are equally detectable:
AUC $0.780$ vs.\ $0.782$ at $\varepsilon{=}0.002$ ($950$ and $650$ examples),
$0.928$ vs.\ $0.939$ at $0.004$ ($1170$ and $430$), and $0.997$ vs.\ $0.999$
at $0.008$. The second group holds all the other images, including the $64$
that the model misclassifies when clean.
% src: experiments/2026-07-21_success-conditioned/outputs/metrics.json [EST — E3]
The probe responds to the perturbation in activation space, not to its effect
on the label.

\subsection{The head split}
\label{sec:adv-head-split}
A displacement can change the label only through the part of it that changes
the logits. The classification head is linear, $Z = Wh + b$ with
$W\in\mathbb{R}^{10\times768}$, on the class token $h$ of the last block after
the final LayerNorm, so the Jacobian of the logits is exact, $J=W$, and each
displacement $\Delta h$ of the head input splits into two orthogonal parts:
\begin{equation}
  \Delta h = \Delta h_{\parallel} + \Delta h_{\perp},
  \qquad \Delta h_{\parallel} = W^{+}W\,\Delta h,
  \label{eq:adv-split}
\end{equation}
where $W^{+}$ is the pseudo-inverse of $W$. The \emph{functional component}
$\Delta h_{\parallel}$ is the only part that changes the logits; the
\emph{non-functional component} satisfies $W\Delta h_{\perp}=0$. The row space
of $W$ has at most $10$ of the $768$ dimensions, so an isotropic displacement
puts about $1.3\%$ of its energy in $\Delta h_{\parallel}$. We train the probe
of Section~\ref{sec:adv-protocol} on $\Delta h$, on $\Delta h_{\parallel}$ and
on $\Delta h_{\perp}$ separately. On FGSM at $\varepsilon{=}0.004$ ($1600$
examples), $87.7\%$ of the energy of $\Delta h$ is in $\Delta h_{\perp}$, and
a probe restricted to $\Delta h_{\perp}$ loses nothing: the AUC is $0.968$ on
the full $\Delta h$, $0.967$ on $\Delta h_{\perp}$ alone, and $0.694$ on
$\Delta h_{\parallel}$ alone, which is above chance.
% src: experiments/2026-07-21_jacobian/outputs/e12a_full/jacobian_metrics.json [EST — E12a]

Three notes limit this result. First, the attack does concentrate energy in
$\Delta h_{\parallel}$: the measured $12.3\%$ is about $10\times$ the isotropic
expectation of $1.3\%$, so $\Delta h_{\parallel}$ is not empty of signal.
Second, the two subspaces have very different dimensions (at most $10$ and
about $758$), so the AUCs show where the signal is and are not a paired
comparison. Third, the split is exact at the head input and not at
\texttt{mlp\_10}, because the map from \texttt{mlp\_10} to the logits is not
linear; there we use the local per-example Jacobian of the logits with respect
to the mean-pooled \texttt{mlp\_10} activation, evaluated at the clean input,
in place of $W$. The AUC of $0.968$ is that of a probe on the head input,
which is lower than that of the probe on $a$ at \texttt{mlp\_10} ($0.989$).
% src: src/adv_interp/cli/compute_jacobian.py:79-81,126-131,195-200; analysis/jspace.py:42-44; EXPERIMENT_PLAN_ATTACKS.md:194-216

One case stays open: an attack that is built to have a small displacement
(Section~\ref{sec:adv-attackers}).

\subsection{The probe does not read only the size of the displacement}
\label{sec:adv-not-size}
The scalar norms ($\lVert r\rVert$, $\lVert\hat{a}\rVert$, $\lVert z\rVert_1$)
lose the detection signal: they detect FGSM at $0.57$--$0.59$, where the
vector representations reach $0.93$--$0.99$
(Table~\ref{tab:adv-representations}). This agrees with a probe that reads the
direction of the displacement and not only its size.
% src: 2026-06-30_decomp-eps0004 decomposition_metrics.csv; 2026-07-02_signal-zoo/outputs/zoo/zoo_metrics.csv [EST]
The budget sweep gives a second example: PGD at $\varepsilon{=}0.002$ and FGSM
at $\varepsilon{=}0.004$ arrive with the same mean displacement ($2.01$ and
$2.00$), and they are detected at different AUCs on $z$ ($0.827$ and $0.932$;
Table~\ref{tab:adv-h1}).
% [EST — the two numbers are in tab:adv-h1 and in sec:adv-amplification; the comparison itself is from ACTION_POINTS_ARTUR.md B1, not from main.tex]

A Mahalanobis distance in the simplest form of \citeonline{lee2018mahalanobis}
(one Gaussian on the clean training-fold activations at one layer), on the
same mean-pooled activations and splits, reaches $0.710$ on FGSM at
$\varepsilon{=}0.004$ and \texttt{mlp\_10}, where the logistic probe reaches
$0.989$, and it is near chance on the input-matched benign families
($0.482$--$0.557$) and on the watermark control ($0.498$).
% src: experiments/2026-07-21_mahalanobis/outputs/maha_vs_logit.csv; maha_baseline.py (single clean Gaussian, shrinkage 1e-3, per-layer) [EST — E8]
The baseline is run without the layer-wise ensemble and the input
pre-processing of the original method, so the claim is not that the method
fails; it is that the probe reads a learned direction in displacement space
(Equation~\eqref{eq:adv-probe}), not a generic off-manifold distance.

\subsection{What remains at equal displacement: the collapse test}
\label{sec:adv-collapse}
\emph{Objection: ``$\dauc$ is negative, not zero, so something remains.''}

What remains is direction. By Equation~\eqref{eq:adv-probe}, the change of
the probe score depends on $\lVert\Delta a\rVert$ and on $\cos\theta$. The
displacement-matched control matches $\lVert\Delta a\rVert$ but not $\theta$,
so families with equal displacement can still get different scores. The
collapse test asks what this direction separates.

We train one clean-vs-FGSM probe on $z$ at \texttt{mlp\_10} (held-out AUC
$0.933$) and score the benign families with it, without training it again. We
pair each score with the $\lVert\Delta a\rVert$ of its example, and a nested
$F$-test asks whether family identity adds explanatory power to a polynomial
in $\lVert\Delta a\rVert$.
% src: experiments/2026-07-22_collapse-e5/outputs/metrics.json (protocol block) [EXPL — E5, expectations pre-stated]
If the probe read only $\lVert\Delta a\rVert$, all families would fall on one
curve. They do not: family identity adds explanatory power
($\Delta R^2=0.373$, $F=1204$, $8000$ examples, $p\approx0$;
Figure~\ref{fig:adv-collapse}). % src: same file (collapse_test) [EST]
In the matched band $\lVert\Delta a\rVert\in[1.8,2.2]$, Gaussian noise scores
as high as FGSM ($0.890$ against $0.877$), and the smooth families score low
(JPEG $0.291$, blur $0.206$; dropout $0.519$).
% src: same file (per_family_score, matched band) [EST]
This agrees with the view that adversarial examples and errors under Gaussian
noise come from one property of the model
\cite{gilmer2019adversarial_examples_noise}.

\begin{figure}[ht]
  \centering
  \includegraphics[width=0.6\linewidth]{figures/interp_adv/fig_collapse.pdf}
  \caption{Collapse test at \texttt{mlp\_10}: mean score of the clean-vs-FGSM
  probe on $z$, applied without retraining to each family, in bins of
  $\lVert\Delta a\rVert$ ($1600$ examples per family, displacement-matched
  families). The shaded band is the matched band
  $\lVert\Delta a\rVert\in[1.8,2.2]$. The families do not fall on one curve: at
  equal displacement, Gaussian noise scores like FGSM, and the smooth families
  score low.}
  % src: figures/src/appendix_items.py; 2026-07-22_collapse-e5/outputs/collapse_table.csv
  \label{fig:adv-collapse}
\end{figure}

The direction that the FGSM probe learns therefore separates noise-like from
smooth perturbations, not attacks from benign perturbations. We call this
property \emph{texture}, as an operational hypothesis. We do not measure the
frequency content of the perturbations, and dropout, whose sparse impulses
carry high-frequency energy, falls between noise and the smooth families. The
test measures how well the FGSM direction transfers to each family, not how
detectable each family is: with its own probe, each smooth family is detected
at least as well as FGSM (Section~\ref{sec:adv-main}).
% NOTE (B13): the research journal treats this single-detector result as a
% candidate for the main result of the paper (partial, three of five seeds).
% If the paper adopts it, this subsection moves up and the per-family probes of
% sec:adv-main become the supporting evidence.

\paragraph*{The refined claim.}
Together with Section~\ref{sec:adv-main}, this result gives the refined form
of $H_{D}$: the probe reads the displacement along a learned direction, and at
equal displacement this direction separates noise-like from smooth
perturbations. Adversarial optimization is a cause of the large displacement,
but the probe does not read adversarial meaning.

% =============================================================================
\section{Step 5: two attackers, a smaller displacement and a different direction}
\label{sec:adv-attackers}
\emph{Objection: ``If the probe reads displacement, an attacker can evade it
with an example that has a low displacement.''}

FGSM, PGD and CW-$L_2$ do not try to keep the displacement small. We therefore
built an attack that does, and a second attack that uses the probe itself.

\subsection{The minimal-displacement attack}
\label{sec:adv-lowdx}
For each image, the attack minimizes the displacement at the probed layer,
\begin{equation}
  \min_{\tilde{x}} \;\lVert a_{10}(\tilde{x}) - a_{10}(x_i)\rVert_2^2
  \;+\; c\cdot\max\bigl(\mathrm{margin}(\tilde{x}), -\kappa\bigr),
  \label{eq:adv-mindisp}
\end{equation}
where $a_{10}$ is the mean-pooled \texttt{mlp\_10} activation and
$\mathrm{margin}(\tilde{x})$ is the logit of the true class minus the largest
other logit. A tanh reparametrization keeps $\tilde{x}$ in pixel space, and an
outer binary search finds the smallest $c$ that still changes the label. The
\emph{floor} is the smallest realized $\lVert\Delta a\rVert$ among successful
iterates; by construction, it is an upper bound on the true minimum. We fixed
the benign band in advance at $\lVert\Delta a\rVert \leq 0.30$, anchored on
the JPEG displacement in the input-matched control ($0.284$).
% src: experiments/2026-07-25_min-displacement-cw/README.md (pre-registration block); params.yaml; src/adv_interp/attacks/cw_min_displacement.py [EXPL — outcome branches pre-stated before the full run]

\paragraph*{The low-displacement cell can be occupied.}
On all $1536$ images that the model classifies correctly when clean, the
attack succeeds (ASR $1.000$), with a median floor of
$\lVert\Delta a\rVert = 0.298$ at \texttt{mlp\_10} (IQR $[0.235, 0.368]$;
$50.7\%$ inside the benign band). This is $6.7\times$ below FGSM at
$\varepsilon{=}0.004$ ($2.00$) and $6.1\times$ below CW-$L_2$ ($1.81$); one
block later, at \texttt{mlp\_11}, the same images have $0.579$.
% src: experiments/2026-07-25_min-displacement-cw/outputs/full/floor_summary.json; verdict.json [EST]
The objective has no constraint on the input norm. The median pixel $L_\infty$
is $0.137$, $34\times$ the FGSM budget at $\varepsilon{=}0.004$; the input
$L_2$ is $8.39$, $1.22\times$ that of FGSM; and the perturbations are visible.
% src: same floor_summary.json (input_linf/input_l2 blocks) [EST; L2 ratio calculated, 8.39/6.85]
For this attack, a small displacement requires a large input perturbation, and
the size of this cost depends on the norm: it is large in $L_\infty$ and small
in $L_2$. The attack does not use the probe: it is adaptive to our hypothesis
about the probe, not to the probe itself.

The search reaches a stable value: the median displacement of the successful
examples stays between $0.308$ and $0.315$ over the first five values of $c$,
and it rises at the smaller values, where fewer attacks succeed. This per-step
median is not the per-image floor (median $0.298$), which is the smallest
successful displacement of each image over all steps. The search used about
half of the planned inner budget, so each floor is an upper bound on the
minimal displacement of its image, and a longer search could find lower values.
% src: outputs/full/pareto_trace.csv; params.yaml (attack.search_budget_rationale) [EST — upper-bound discipline]

\paragraph*{The probe is not evaded.}
The AUC of clean against the minimal-displacement examples is $0.999$ on $a$
and $0.990$ on $z$ at \texttt{mlp\_10}.
% src: outputs/full/detection_auc.json (floor_set cells) [EST]
At almost the same displacement, JPEG in the input-matched control
($\lVert\Delta a\rVert = 0.284$) is detected at $0.604$ on $z$, which is an
asymmetry of $\approx 0.39$ AUC at matched displacement.
% src: detection_auc.json; 2026-06-29_h2-v2-eps0004/outputs/h2_crossval/metrics.json (benign_jpeg) [EST]
\pendingpaper{A4}{We report this asymmetry without an interpretation, because
one alternative explanation stays open. The input perturbation of these
examples is much larger than that of JPEG, so its small displacement can point
in a direction that a linear probe separates easily ($\cos\theta$ in
Equation~\eqref{eq:adv-probe}). A benign family matched to their input
$L_\infty$ would separate ``adversarial structure beyond displacement'' from
``a direction caused by a large input perturbation''; this control was not
run. This cell is the one open question of the chapter.}
% [DECISION 2026-07-28 (Breno): reported, not interpreted; L_inf-matched control not run]

The head split on these examples agrees with
Section~\ref{sec:adv-head-split}. $92.4\%$ of the energy of their displacement
is in $\Delta h_{\perp}$ (FGSM: $87.7\%$), and the probe signal is there: AUC
$0.9996$ on the full $\Delta h$, $0.9997$ on $\Delta h_{\perp}$ alone, and
$0.7804$ on $\Delta h_{\parallel}$ alone. At \texttt{mlp\_10}, with the local
Jacobian, the functional component alone is at chance ($0.518$).
% src: outputs/full/perp_parallel.json [EST — E12-consistent]

\subsection{An adaptive attack against the probe}
\label{sec:adv-adaptive}
The minimal-displacement attack does not use the probe, so we also ran an
attack that does, with a decision rule pre-specified before the run.
% src: JOURNAL_V3.md §8c (P8, frozen 2026-10-27, with pre-results amendments) [PRE — P8]
The threat model is white-box on the classifier and on the probe: the attacker
knows the ViT-B/16, the pooling, the standardization and the weights of the
logistic probe. The attacked probe is the clean-vs-APGD-CE probe at
\texttt{mlp\_10}, one on $a$ and one on $z$; each image is attacked with the
probe of the cross-validation fold that holds it out. For each fold, the
threshold $\tau$ is the $95$th percentile of the probe score on the held-out
clean images (false-positive rate $5\%$). The attack is successful on an image
if the label is wrong and the probe score is below $\tau$. The adaptive APGD
(Table~\ref{tab:adv-attacks}) maximizes
\begin{equation}
  L(\tilde{x}) \;=\; \mathrm{CE}\bigl(g(\tilde{x}), y\bigr) \;-\; \lambda\, s(\tilde{x}),
  \label{eq:adv-adaptive}
\end{equation}
where $g$ is the classifier, $y$ the true label and $s$ the probe logit, at
$\varepsilon{=}0.004$. The attack uses $100$ iterations and one restart, with
$\lambda\in\{0.1, 1, 10\}$ in sequence: a larger $\lambda$ is applied only to
the images on which the smaller values failed. The decision rule defines
evasion as a success rate of at least $0.5$ on the images that the model
classifies correctly when clean.

\paragraph*{The attacker evades the probe.}
Over five seeds, the success rate is $0.774\pm0.010$ on $a$ and
$0.730\pm0.017$ on $z$, and almost all successes occur at $\lambda=0.1$.
% src: experiments/2026-10-27_p8-adaptive/outputs/metrics.json [PRE — P8 gate: evasion in both cells]
No final example has a probe score above $\tau$: the attack fails only on
images whose label does not change.
% src: outputs/{pilot,runs}/seed*/cell_*_eps0.004/per_example.csv (s_adv vs tau, pred_adv vs label; computed for the paper) [EST — not in the journal]
Thus, for this attacker the probe constraint is easy and the label constraint
is the difficult one. The success also shows directly that the gradient of the
probe is not masked. The input cost is that of FGSM, an $L_\infty$ of $0.004$.

\paragraph*{The attacker does not remove the displacement.}
The attacker evades the probe by changing the direction of the displacement,
not its size (Figure~\ref{fig:adv-adaptive}). The median
$\lVert\Delta a\rVert$ at \texttt{mlp\_10} is $2.55$ on $a$ and $2.82$ on $z$,
against $2.38$ for APGD-CE on the same images, and more than $99.9\%$ of the
adaptive examples are above the largest benign displacement in the
input-matched control ($0.59$). What changes is $\cos\theta$ in
Equation~\eqref{eq:adv-probe}, measured between the change of the probed
representation and the probe direction: its median goes from $+0.092$ for
APGD-CE to $-0.032$ for the adaptive attack on $a$. On $z$ the values are much
smaller, and the sign changes in the same way (from $+0.0003$ to $-0.0023$).
% src: experiments/2026-10-27_p8-adaptive/outputs/diag/diag_metrics_with_retrained.json; diag_per_example.csv [PRE diagnostic, reported not gated]
A probe retrained on clean against adaptive examples, with the same protocol,
detects them at AUC $0.991$ on $a$ and $0.998$ on $z$. The clean-vs-FGSM
probe, which the attacker did not use, gives the adaptive examples lower
scores than the clean images on $a$ (AUC $0.34$--$0.44$, three seeds) and is
near chance on $z$ ($0.51$--$0.66$); the attack moved the examples to the
clean side of the direction that the two probes share.
% src: diag_metrics_with_retrained.json (retrained_probe, transfer_p1; seeds 42/44/46)

\begin{figure}[ht]
  \centering
  \includegraphics[width=\linewidth]{figures/interp_adv/fig_adaptive.pdf}
  \caption{The adaptive attack changes the direction of the displacement, not
  its size. Each point is one image that the model classifies correctly when
  clean (five seeds): $\lVert\Delta a\rVert$ at \texttt{mlp\_10} against
  $\cos\theta$ between the change of the probed representation and the probe
  direction, for APGD-CE and for the adaptive APGD at $\varepsilon{=}0.004$.
  White circles are medians. The dashed line is the largest benign
  displacement in the input-matched control. Adaptive examples that keep the
  true label are in grey; all adaptive examples are below $\tau$. The $z$
  panel does not show the $1\%$ most negative adaptive values.}
  % src: figures/src/fig_adaptive.py (paper repository)
  \label{fig:adv-adaptive}
\end{figure}

\subsection{What the two attackers show}
An attacker who knows the probe evades it by direction at the FGSM budget, and
the displacement stays at the attack scale; a retrained probe detects these
examples again. Only the costly minimal-displacement attack removes the
displacement, and the probe still detects it. The two attackers are the two
limits of $H_{D}$ stated in Section~\ref{sec:adv-two-hypotheses}: a large
displacement is sufficient for detection by a probe trained on that family,
not by a fixed probe against an attacker who knows it; and a small
displacement does not imply that the example is not detected.

% =============================================================================
\section{Layers and representations}
\label{sec:adv-layers}
% Thesis-only material (not in the paper, or reduced there to one sentence).

The five steps use one probed layer, \texttt{mlp\_10}, and one budget,
$\varepsilon{=}0.004$. This section gives the sweeps over layers and
representations that support these two choices.

\subsection{The saturated budget}
\label{sec:adv-saturated}
At $\varepsilon{=}0.012$, the budget of the pre-specified rules
(Section~\ref{sec:adv-provenance}), every representation detects FGSM at the
ceiling at every probed layer (Table~\ref{tab:adv-decomp},
Figure~\ref{fig:adv-decomp-layers}a). At the ceiling, a comparison between
representations or between families says nothing, and this is the reason for
the range before saturation that the five steps use. The watermark control is
near chance from \texttt{mlp\_6} on; at the first two blocks it reaches
$0.675$--$0.685$ on $a$, and it decays with depth.
% src: experiments/2026-06-22_sae-residual-decomp/outputs/decomposition_metrics.csv [EST]

\begin{table}[ht]
  \centering
  \caption{Detection AUC (grouped cross-validation, $K{=}5$, $N{=}3200$) by
  representation and layer at the \emph{saturated} budget
  ($\varepsilon{=}0.012$). FGSM is at the ceiling everywhere; the watermark
  control is near chance at the late layers.}
  \label{tab:adv-decomp}
  \small
  \begin{tabular}{@{}llcccc@{}}
    \toprule
    Layer & Family & $a$ & $z$ & $\hat{a}$ & $r$ \\
    \midrule
    \multirow{2}{*}{\texttt{mlp\_0}}  & FGSM      & 1.000 & 0.996 & 1.000 & 1.000 \\
                                      & Watermark & 0.675 & 0.539 & 0.541 & 0.677 \\
    \multirow{2}{*}{\texttt{mlp\_3}}  & FGSM      & 1.000 & 1.000 & 1.000 & 1.000 \\
                                      & Watermark & 0.685 & 0.604 & 0.604 & 0.664 \\
    \multirow{2}{*}{\texttt{mlp\_6}}  & FGSM      & 1.000 & 1.000 & 1.000 & 1.000 \\
                                      & Watermark & 0.549 & 0.531 & 0.537 & 0.542 \\
    \multirow{2}{*}{\texttt{mlp\_9}}  & FGSM      & 1.000 & 0.999 & 0.998 & 1.000 \\
                                      & Watermark & 0.530 & 0.513 & 0.524 & 0.530 \\
    \multirow{2}{*}{\texttt{mlp\_10}} & FGSM      & 1.000 & 1.000 & 1.000 & 1.000 \\
                                      & Watermark & 0.543 & 0.519 & 0.532 & 0.540 \\
    \multirow{2}{*}{\texttt{mlp\_11}} & FGSM      & 1.000 & 0.997 & 1.000 & 1.000 \\
                                      & Watermark & 0.544 & 0.513 & 0.524 & 0.542 \\
    \bottomrule
  \end{tabular}
  % src: experiments/2026-06-22_sae-residual-decomp/outputs/decomposition_metrics.csv
\end{table}

\begin{figure}[ht]
  \centering
  \includegraphics[width=\linewidth]{figures/interp_adv/fig_decomp_layers.pdf}
  \caption{Detection across layers. (a) At the saturated budget
  ($\varepsilon{=}0.012$) every representation detects FGSM at the ceiling at
  every block, while the watermark control decays to chance. (b) In the
  input-matched control at $\varepsilon{=}0.004$, the detection advantage
  $\dauc$ is negative or near zero at the first blocks and largest at
  \texttt{mlp\_10}. In the legend, $x$ denotes the raw activation $a$ of
  Equation~\eqref{eq:adv-decomp}, and ``stego'' the watermark control.}
  % src: experiments/2026-06-22_sae-residual-decomp/outputs/decomposition_metrics.csv;
  %      experiments/2026-07-21_decomp-6layer/outputs/decomposition_metrics.csv;
  %      experiments/2026-07-02_gap-robustness/outputs/gap_robustness_summary.csv
  % The figure predates the notation change (x -> a) and the rename of the
  % control; regenerate it in the research repository when possible.
  \label{fig:adv-decomp-layers}
\end{figure}

\subsection{The input-matched advantage across layers}
\label{sec:adv-depth}
The detection advantage in the input-matched control is a property of the mid
and late layers (Figure~\ref{fig:adv-decomp-layers}b): it is negative at
\texttt{mlp\_0} and rises to $0.204\pm0.014$ on $z$ at \texttt{mlp\_10} over
five seeds, while FGSM is already detected at the first block. The curve on
$a$ comes from the six-layer decomposition, which ran at a reduced sample
($192$ examples). A possible explanation is the amplification of
Section~\ref{sec:adv-amplification}: the displacement of FGSM grows with depth,
from $0.09$ at \texttt{mlp\_0} to $2.00$ at \texttt{mlp\_10}
(Table~\ref{tab:adv-dx-layers}). We did not measure the displacement of the
input-matched benign families at the layers other than \texttt{mlp\_10}, so
this explanation is a conjecture.
% src: experiments/2026-07-02_gap-robustness/outputs/gap_robustness_summary.csv; 2026-07-21_decomp-6layer [EST]

\subsection{The displacement-matched control on the latents, per layer}
\label{sec:adv-matched-dx-z}
Table~\ref{tab:adv-dx-layers} gives $\dauc$ on $a$ at each layer in the
displacement-matched control. Table~\ref{tab:adv-matched-dx} gives the same
control on $z$, with the FVE of the SAE at each layer. Only \texttt{mlp\_10}
and \texttt{mlp\_11} pass the health gate, and the caveat of
Section~\ref{sec:adv-dx-layers} applies to both tables: the match is
calibrated at \texttt{mlp\_10}.
% src: experiments/2026-07-22_h2-matched-dx/outputs/h2_matched_dx_summary.json [EST — E1]

\begin{table}[ht]
  \centering
  \caption{Displacement-matched control on the latents $z$, seed 42.
  $^{\dagger}$FVE below the $0.9$ health gate: these cells are reported for
  completeness only. Five-seed replication: $\dauc = -0.063\pm0.003$ at
  \texttt{mlp\_10} and $-0.076\pm0.005$ at \texttt{mlp\_11}.}
  \label{tab:adv-matched-dx}
  \small
  \begin{tabular}{@{}lcccc@{}}
    \toprule
    Layer & FVE & FGSM, AUC on $z$ & Best benign, AUC on $z$ & $\dauc$ on $z$ \\
    \midrule
    \texttt{mlp\_0}  & $0.753^{\dagger}$ & $0.944$ & $0.999$ & $-0.055$ \\
    \texttt{mlp\_3}  & $0.779^{\dagger}$ & $1.000$ & $1.000$ & $-0.000$ \\
    \texttt{mlp\_6}  & $0.720^{\dagger}$ & $0.977$ & $1.000$ & $-0.023$ \\
    \texttt{mlp\_9}  & $0.697^{\dagger}$ & $0.886$ & $0.994$ & $-0.108$ \\
    \texttt{mlp\_10} & $0.936$ & $0.933$ & $0.995$ & $-0.062$ \\
    \texttt{mlp\_11} & $0.947$ & $0.918$ & $0.996$ & $-0.078$ \\
    \bottomrule
  \end{tabular}
  % src: experiments/2026-07-22_h2-matched-dx/outputs/h2_matched_dx_summary.json;
  %      experiments/2026-07-23_h2-matched-dx-multiseed/outputs/multiseed_summary.csv
\end{table}

\subsection{Eleven detection signals}
\label{sec:adv-zoo}
The SAE is a lossy bottleneck, so $z$ cannot carry detection information that
$a$ lacks; a linear probe on $z$ could still gain from the non-linearity of
the encoder, but here it does not (Section~\ref{sec:adv-scope}). A sweep of
eleven detection signals in the same cell (input-matched control,
$\varepsilon{=}0.004$, \texttt{mlp\_10}; Figure~\ref{fig:adv-zoo}) orders them
by $\dauc$: the signals on the SAE features have the largest input-matched
$\dauc$, and the scalar magnitudes the smallest. The anti-condition of
Section~\ref{sec:adv-scope} applies to this ordering too: it compares
differences of AUCs at different distances from the ceiling.
% src: experiments/2026-07-02_signal-zoo/outputs/zoo/zoo_metrics.csv [EST]

\begin{figure}[ht]
  \centering
  \includegraphics[width=0.85\linewidth]{figures/interp_adv/fig_zoo.pdf}
  \caption{Input-matched $\dauc$ (FGSM minus the most detectable input-matched
  benign family) for eleven detection signals at $\varepsilon{=}0.004$,
  \texttt{mlp\_10}. In the labels, $x$ and $\hat{x}$ denote the raw activation
  $a$ and its reconstruction $\hat{a}$ of Equation~\eqref{eq:adv-decomp}.}
  % src: experiments/2026-07-02_signal-zoo/outputs/zoo/zoo_metrics.csv
  % The figure predates the notation change (x -> a).
  \label{fig:adv-zoo}
\end{figure}

% =============================================================================
\section{Provenance of the tests}
\label{sec:adv-provenance}
% The paper's provenance appendix, plus the frozen-budget runs of the two
% pre-specified rules, which only the thesis reports in full.

The order of execution was different from the order of
Sections~\ref{sec:adv-usual} to~\ref{sec:adv-attackers}. A test is
\emph{pre-specified} when its hypothesis and decision rule were committed to
the version-controlled project repository before the run; this is not a formal
registration with a third party. Before any run, we pre-specified two
predictions of $H_{M}$ with fixed decision rules; the rules are never
evaluated again on later data.
% src: EXPERIMENT_PLAN_ATTACKS.md (rev. 2026-07-21, rule block); PLAN.md; DESIGN_CHOICES.md [PRE]

\paragraph*{The two rules.}
\emph{H1 (budget invariance)} passes if, over $\varepsilon\leq0.012$, the
correlation between $\varepsilon$ and the AUC is near zero ($|\rho|<0.3$), the
variance of the AUC is below $10^{-3}$, and the mean AUC is above $0.95$.
\emph{H2 (adversarial specificity)} passes if, at the frozen budget
$\varepsilon{=}0.012$ and against benign families of matched nominal $L_2$,
$\dauc>0.3$ and the most detectable benign family is below AUC $0.7$.
% src: PLAN.md; DESIGN_CHOICES.md (frozen-gate definitions) [PRE]

\paragraph*{H1 fails.}
H1 fails all three of its criteria for FGSM and PGD, because the AUC rises
monotonically with the budget: from $0.670$ to $1.000$ for FGSM, with
saturation at $\varepsilon\geq0.012$, and from $0.663$ to $1.000$ for PGD
(Table~\ref{tab:adv-h1}, Figure~\ref{fig:adv-h1}). The mean over the five
budgets $\varepsilon\leq0.012$ is about $0.88$ for FGSM, below the floor.
% src: both magnitude_sweep_metrics.csv; RESEARCH_LOG.md (corrected re-run verdict) [PRE — gate fails]
% The 0.88 is arithmetic on the five FGSM values (0.670, 0.785, 0.932, 0.996, 1.000).

\paragraph*{H2 fails, at the frozen budget and below it.}
At the frozen budget the attack is at the ceiling (FGSM $0.9999$), and every
benign family of comparable nominal magnitude is detected at $0.85$--$0.94$
(JPEG $0.939$, dropout $0.921$, blur $0.910$, noise $0.847$). $\dauc = 0.061$,
and the rule fails on both criteria.
% src: experiments/2026-06-23_corrected-rerun/outputs/h2_crossval/metrics.json [EST/PRE]
This run is older than the per-family measurement of the input $L_2$, so at
the frozen budget the input match is the design intent of the perturbation
configurations, not a verified match. In the range before saturation, with the
match verified, the rule still fails: $\dauc = 0.199$ at $\varepsilon{=}0.004$
is below $0.3$, and the most detectable benign family (blur, $0.733$) is above
$0.7$ (Section~\ref{sec:adv-usual}). The input-matched $\dauc$ also passes a third pre-specified rule, a robustness criterion
($\dauc-2\sigma > 0.05$ at $\varepsilon\in\{0.002,0.004\}$, at mid and late
layers, over five seeds); this rule tests whether a detection advantage
exists, and H2 tests whether it is large.
% src: gap_robustness_summary.csv; DESIGN_CHOICES.md D12 (G5.3, frozen 2026-07-02) [PRE — G5.3 passes]

\begin{table}[ht]
  \centering
  \caption{H1 budget sweep: detection AUC of the clean-vs-attack probe on $z$
  at \texttt{mlp\_10} (grouped 5-fold cross-validation, $N{=}3200$ per task)
  and attack success rate (ASR), for FGSM and PGD. The AUC rises with
  $\varepsilon$ and saturates (FGSM at $\varepsilon\geq0.012$, PGD at
  $\varepsilon\geq0.016$); H1 predicted a constant AUC.}
  \label{tab:adv-h1}
  \small
  \begin{tabular}{@{}ccccc@{}}
    \toprule
     & \multicolumn{2}{c}{FGSM} & \multicolumn{2}{c}{PGD} \\
    \cmidrule(lr){2-3}\cmidrule(lr){4-5}
    $\varepsilon$ & AUC & ASR & AUC & ASR \\
    \midrule
    $0.001$ & $0.670$ & $0.382$ & $0.663$ & $0.580$ \\
    $0.002$ & $0.785$ & $0.619$ & $0.827$ & $0.922$ \\
    $0.004$ & $0.932$ & $0.762$ & $0.969$ & $0.999$ \\
    $0.008$ & $0.996$ & $0.819$ & $0.993$ & $1.000$ \\
    $0.012$ & $1.000$ & $0.832$ & $0.999$ & $1.000$ \\
    $0.016$ & $1.000$ & $0.831$ & $1.000$ & $1.000$ \\
    $0.024$ & $1.000$ & $0.817$ & $1.000$ & $1.000$ \\
    $0.032$ & $1.000$ & $0.808$ & $1.000$ & $1.000$ \\
    \bottomrule
  \end{tabular}
  % src: experiments/2026-06-23_corrected-rerun/outputs/h1_sweep/magnitude_sweep_metrics.csv;
  %      experiments/2026-07-02_pgd-h1h2/outputs/h1_sweep/magnitude_sweep_metrics.csv
\end{table}

\begin{figure}[ht]
  \centering
  \includegraphics[width=0.75\linewidth]{figures/interp_adv/fig_h1_sweep.pdf}
  \caption{Budget sweep: detection AUC against the budget for FGSM and PGD,
  with the AUC of successful and of failed FGSM attacks
  (Section~\ref{sec:adv-success}). The AUC rises with the budget, and
  successful and failed attacks are equally detectable.}
  % src: the two magnitude_sweep_metrics.csv; experiments/2026-07-21_success-conditioned/outputs/metrics.json
  \label{fig:adv-h1}
\end{figure}

\paragraph*{What the failures show.}
The two rules fail, for two attacks. This eliminates $H_{M}$ in the form that
was pre-specified, and it is weak evidence for $H_{D}$, for two reasons. The
first is saturation: when the AUC is at the ceiling, H1 cannot separate a
detector of displacement from a detector of adversarial structure that grows
with $\varepsilon$, and the $\dauc$ of H2 is measured against the ceiling. The
success split of Section~\ref{sec:adv-success} answers the first reading, and
the runs at $\varepsilon\approx0.002$--$0.004$ answer the second. The second
reason is logical. Because the two rules test $H_{M}$, their failure is only
consistent with $H_{D}$ and does not test it. A test counts as evidence for $H_{D}$ only if it
could have failed under $H_{D}$ \cite{mayo2018severe_testing}; this is the
method of strong inference with multiple working hypotheses
\cite{platt1964strong_inference,chamberlin1890multiple_working_hypotheses}.
For this reason, the failures of H1 and of H2 carry less weight than the
displacement-matched control (Section~\ref{sec:adv-main}). The later
experiments test predictions of $H_{D}$ itself: the displacement-matched
control (Section~\ref{sec:adv-main}), the success split and the head split
(Sections~\ref{sec:adv-success} and~\ref{sec:adv-head-split}), and the
collapse test (Section~\ref{sec:adv-collapse}). These experiments are
exploratory: we wrote down their expected outcomes before execution, but they
were not pre-specified.
% src: EXPERIMENT_PLAN_ATTACKS.md (rev. 2026-07-21) [EXPL]
The two APGD experiments came last, and both were pre-specified with fixed
decision rules: the APGD-CE replication (Section~\ref{sec:adv-apgd-ce}) and
the adaptive attack against the probe (Section~\ref{sec:adv-adaptive}).
% src: JOURNAL_V3.md §8b (P7, frozen 2026-10-19), §8c (P8, frozen 2026-10-27) [PRE]

% =============================================================================
\section{Two measurement bugs}
\label{sec:adv-bugs}
% Thesis-only material. Kept because each bug gave a plausible result that was
% wrong, and because the first one motivated the statistical controls of
% sec:adv-protocol.

Two measurement bugs were found during this work. Each gave a plausible result
that was wrong, and each gives a rule that applies to other studies.

\paragraph*{The clip bug: an artefact that looked like budget-invariant
detection.}
The FGSM routine first applied $\mathrm{clamp}(\cdot, 0, 1)$ to
ImageNet-\emph{normalized} images, whose values span about $[-2.2, 2.6]$.
% src: DESIGN_CHOICES.md; RESEARCH_LOG.md 2026-04-23 [EST]
Clipping in that space changes most pixels: at $\varepsilon{=}0.012$ the
measured per-image $L_2$ was about $284$, against the $20$ that the budget
implies, and the clipped images were far from every clean image, in a way that
did not depend on $\varepsilon$. The detector read this artefact. The sweep
before the fix was therefore \emph{flat} at AUC $\approx 0.88$, a false
confirmation of H1: the artefact added about $0.23$ AUC at
$\varepsilon{=}0.001$ and only $0.02$ at $0.032$, and it cancelled the true
trend with the budget.
% src: RESEARCH_LOG.md 2026-04-23 (delta table) [EST]
The fix clips in pixel space (denormalize, clip in $[0,1]$, renormalize);
after the fix, the $L_2$ at $\varepsilon{=}0.012$ is $20.51$, which agrees
with the budget, and the sweep gives the monotone curve of
Table~\ref{tab:adv-h1}.
% src: src/adv_interp/attacks/fgsm.py [EST]
The rule: measure the realized perturbation norm and compare it with the
budget \emph{before} reporting a detection number. With the bug, the flat
sweep would have passed the two invariance criteria of the pre-specified H1
rule (correlation and variance) on an artefact, and failed only the criterion
on the mean AUC.

\paragraph*{The first-key bug: the SAE was trained on a different layer from
the documented one.}
The SAE trainer took as its target the first non-label key of the saved
activation file. With forward hooks registered in the order
$[\texttt{layer\_10}, \texttt{layer\_11}]$, that key was always
\texttt{layer\_10}, and the \texttt{target\_layer} field of the configuration
files had no effect.
% src: CLUSTER_NOTES.md 2026-06-22; RESEARCH_LOG.md [EST]
The reference detection number of the project was thus computed on
\texttt{mlp\_10} alone, not on \texttt{mlp\_11} or on the two layers together,
as the documentation said. This is an error of provenance and does not
invalidate the measurement: the layer sweep, which uses one layer per run, was
not affected. The fix passes an explicit \texttt{--target-layer} argument to
SAE training and to latent encoding; \texttt{mlp\_10} was then adopted as the
probed layer of the main result.
% src: CLUSTER_NOTES.md (P5 rewire) [EST]
The rule: an analysis target must be an explicit, logged argument, never the
result of the order of a container.

% =============================================================================
\section{Reconciliation with deployment work}
\label{sec:adv-reconciliation}
% All three agreements verified against arXiv:2605.07447 (HTML) on 2026-07-24;
% paraphrase only.

SAEgis \cite{wang2026saegis} is an SAE-based adversarial detector for
vision--language models, built for deployment. Where the two studies overlap,
they agree. Its dense baseline without an SAE ties the SAE detector in-domain,
and beats it on one domain, which agrees with the comparison of $a$ and $z$ in
Section~\ref{sec:adv-scope}. Its SAE reconstruction error, which is the scalar
$\lVert r\rVert$ of this chapter up to a monotone transform (the per-example
reconstruction error is $\lVert r\rVert^2/d_{\mathrm{model}}$), shows a
negligible clean--adversarial difference, which agrees with the input-matched
$\dauc$ of $0.054$ on $\lVert r\rVert$
(Table~\ref{tab:adv-representations}). Its most robust cross-attack detector
is at the earliest vision block, and its authors attribute it to
high-frequency texture, which agrees with the noise-like direction of
Section~\ref{sec:adv-collapse}.
% [EXT — all three verified 2026-07-24: (1) dense baseline = cosine to averaged embeddings on raw hidden states; ties on two domains and beats SAEgis on the medical domain; (2) reconstruction-MSE anomaly score shows negligible clean-vs-adversarial difference; (3) vision block 0 most robust cross-attack, high-frequency-texture hypothesis]
We contest only the strong reading, that such detectors read adversarial-class
meaning; the weak reading, that the features give a useful detector, stands.

% =============================================================================
\section{Limitations and the reading for the thesis}
\label{sec:adv-closing}

\subsection{Limitations}

\paragraph*{The main result is exploratory.}
The decisive control, the displacement-matched one, is exploratory: we wrote
down its expected outcome before the run, and it replicates over five seeds,
but it was not pre-specified. \pendingpaper{A1}{Its pre-specified APGD-CE
replication is at the ceiling of the AUC (Section~\ref{sec:adv-apgd-ce}), so
the next step is to pre-specify it below the ceiling and run it on new seeds.}

\paragraph*{Limits of the method.}
The method excludes only the rival explanations that we state, and this set is
not complete \cite{odonohue2001weaknesses_strong_inference}. Each exclusion
also depends on auxiliary assumptions, for example that a match of the mean
displacement is sufficient. The order of the steps is a reconstruction: most
tests were designed after earlier results were known
\cite{gelman2014statistical_crisis}, and for this reason we separate them from
the pre-specified tests
\cite{nosek2018preregistration_revolution,wagenmakers2012confirmatory_agenda}.
A claim that the attack has no detection advantage is an equivalence claim
\cite{lakens2017equivalence_tests}. Our decision rule (mean $\pm2\sigma$ over
five seeds against a margin of $0.05$ in AUC) is not a formal equivalence
test, and the margin, set before the runs from the variation between seeds of
an earlier multi-seed run (standard deviation of $\dauc$ $\approx0.02$), is a
choice.
% src: interp_adversarial_image/DESIGN_CHOICES.md D12 (G5.3 rationale)

\paragraph*{Scope of the instrument.}
Every absence claim is scoped to a linear probe on mean-pooled activations,
and pooling can hide localized signal (Section~\ref{sec:adv-scope}). The
chapter makes no claim about non-linear or token-level probes.
\pendingpaper{A3}{It also does not test whether a probe trained to separate
attacks from benign perturbations at equal displacement succeeds; ``not
meaning'' is tested here as the change of the label, through the success split
and the head split.}

\paragraph*{Scope of the runs.}
The displacement-matched control is replicated over five seeds, at six layers
on $a$ and at the two layers where the SAE passes its health gate on $z$; it
is calibrated at \texttt{mlp\_10} and matches the mean displacement over
images, not the displacement of each image. The collapse test, CW, the
minimal-displacement attack, the token-level probe and the head split are
single-seed; the adaptive attack uses one budget, and its transfer to the
clean-vs-FGSM probe is measured in three of five seeds. The SAE health gate
passes only at \texttt{mlp\_10} and \texttt{mlp\_11}, the per-layer FVE
disagrees between runs at early layers, and the SAE is trained on the clean
images of the detection tasks (Section~\ref{sec:adv-sae}). We use one dataset
(CIFAR-10, upsampled, over an ImageNet-pretrained backbone), one ViT
classifier and one SAE architecture, and make no claim about vision--language
models.
% src: token_probe_metrics.json; zoo_metrics.csv; 2026-07-23 multiseed; 2026-07-22_collapse-e5; 2026-07-21_cw-attack; 2026-07-25_min-displacement-cw; 2026-07-21_stego-token; 2026-07-21_jacobian; 2026-07-21_decomp-6layer; h2_matched_dx_summary.json [EST — scope]

\subsection{The reading for the thesis}
% HOLD (decision of 2026-11-16): the chapter title and the one-sentence thesis
% reading wait for action point A3. This subsection states only what the
% current evidence supports. THESIS-STATEMENT.md still says "mede-se
% ||Delta x|| e prevê-se a AUC"; the data of this chapter contradict that
% sentence (see the second paragraph below), and it is to be corrected in
% Phase 4.

The chapter asked what an activation-space adversarial detector in a ViT
classifier measures. The usual control, which matches benign perturbations to
the attack in input norm, reports that the attack is more detectable than
every benign perturbation ($\dauc=0.199$); when the benign perturbations are
matched in displacement, this advantage disappears. The explanation that the
detector reads an adversarial category gives way to a geometric one: the probe
reads the displacement of the activation along a direction that it learns.
The geometric quantity is measurable for each example, and it can be
controlled: we matched it, and the advantage disappeared; an attacker rotated
it away from the probe direction, and the detection disappeared at an
unchanged displacement.

The geometric quantity is a vector, not a scalar. The norm
$\lVert\Delta a\rVert$ alone does not predict the detection AUC: the scalar
norms are near chance, an attack at the benign scale of displacement is
detected at $0.999$, an adaptive attack at the attack scale is not detected by
the probe that it knows, and two attacks with equal mean displacement are
detected at different AUCs. What the evidence supports is narrower. A large
displacement is sufficient for detection by a probe trained on that family,
and at equal displacement the attack has no detection advantage over benign
perturbations.

Three things follow. First, the sparse autoencoder is not needed for this
detection: the raw activation detects the attack better than the latents, and
the result at equal displacement holds on both. The SAE is a legitimate
measuring instrument, but it is not the source of the signal. Second, the result was visible only with a control that
the usual evaluations do not have and that the pre-specified hypotheses did
not anticipate; the pre-specified tests kept the negative result honest, and
an exploratory experiment found the explanation, so the two are reported with
different status (Section~\ref{sec:adv-provenance}). Third, we expect the
methodological rule to apply beyond this system: before reading meaning into
an activation-space detector, measure the per-example displacement that it
reads, and match the controls there; a match in input norm is not sufficient.
One question stays open, and it bounds the reading: whether the probe reads
adversarial structure in the minimal-displacement examples
(Section~\ref{sec:adv-lowdx}).

\end{document}
